% ============================================================
%  Проект ИТ-архитектуры ГК «Самолет» — Отчёт
%  Компиляция: xelatex main.tex (2 прохода)
% ============================================================
\documentclass[12pt,a4paper]{article}

% --- Шрифты и язык ---
\usepackage{fontspec}
\usepackage{polyglossia}
\setdefaultlanguage{russian}
\setotherlanguage{english}
\defaultfontfeatures{Ligatures=TeX}
\setmainfont{PT Serif}
\setsansfont{PT Sans}
\setmonofont{PT Mono}
\newfontfamily\cyrillicfont{PT Serif}
\newfontfamily\cyrillicfontsf{PT Sans}
\newfontfamily\cyrillicfonttt{PT Mono}

% --- Межстрочный интервал ---
\usepackage{setspace}
\onehalfspacing

% --- Геометрия страницы ---
\usepackage[left=25mm, right=15mm, top=20mm, bottom=20mm, headheight=14pt]{geometry}

% --- Графика и диаграммы ---
\usepackage{tikz}
\usetikzlibrary{
  shapes.geometric, shapes.multipart, shapes.arrows,
  arrows.meta, positioning, calc, fit, backgrounds,
  decorations.pathreplacing, decorations.markings,
  matrix, trees, patterns, shadows.blur
}
\usepackage{pgfplots}
\pgfplotsset{compat=1.18}
\usepackage{pgfgantt}

% --- Таблицы ---
\usepackage{longtable}
\usepackage{booktabs}
\usepackage{tabularx}
\usepackage{array}
\usepackage{multirow}
\usepackage{colortbl}

% --- Оформление ---
\usepackage{fancyhdr}
\usepackage{titlesec}
\usepackage{enumitem}
\usepackage{caption}
\usepackage{float}
\usepackage{xcolor}
\usepackage{hyperref}
\usepackage{indentfirst}

% --- Цветовая палитра ---
\definecolor{samoletblue}{RGB}{0,102,178}
\definecolor{samoletdark}{RGB}{30,55,90}
\definecolor{samoletlight}{RGB}{200,220,245}
\definecolor{accentgreen}{RGB}{46,139,87}
\definecolor{accentorange}{RGB}{230,126,34}
\definecolor{accentred}{RGB}{192,57,43}
\definecolor{accentteal}{RGB}{22,160,133}
\definecolor{accentpurple}{RGB}{142,68,173}
\definecolor{lightgray}{RGB}{240,240,240}
\definecolor{bpmnblue}{RGB}{210,230,250}
\definecolor{bpmnyellow}{RGB}{255,243,200}
\definecolor{bpmngreen}{RGB}{210,245,210}
\definecolor{bpmnred}{RGB}{250,215,215}
\definecolor{layerbiz}{RGB}{255,245,200}
\definecolor{layerapp}{RGB}{200,225,255}
\definecolor{layerdata}{RGB}{210,250,210}
\definecolor{layerinfra}{RGB}{230,230,230}

% --- Стили заголовков ---
\titleformat{\section}{\Large\bfseries\color{samoletblue}}{\thesection.}{0.5em}{}[\vspace{2pt}\titlerule]
\titleformat{\subsection}{\large\bfseries\color{samoletdark}}{\thesubsection.}{0.5em}{}
\titleformat{\subsubsection}{\normalsize\bfseries}{\thesubsubsection.}{0.5em}{}

% --- Колонтитулы ---
\pagestyle{fancy}
\fancyhf{}
\fancyhead[L]{\small\itshape ГК «Самолет» --- Проект ИТ-архитектуры}
\fancyhead[R]{\small\thepage}
\renewcommand{\headrulewidth}{0.4pt}
\renewcommand{\footrulewidth}{0pt}

% --- Гиперссылки ---
\hypersetup{
  colorlinks=true,
  linkcolor=samoletblue,
  urlcolor=samoletblue,
  citecolor=samoletblue,
  bookmarks=true
}

% --- Общие TikZ-стили ---
\tikzset{
  base node/.style={
    rectangle, rounded corners=3pt, draw=gray!60,
    minimum height=10mm, text width=28mm, align=center,
    font=\small\sffamily, inner sep=4pt
  },
  bpmn start/.style={circle, draw=accentgreen, fill=bpmngreen, line width=1.2pt, minimum size=9mm},
  bpmn end/.style={circle, draw=accentred, fill=bpmnred, line width=2pt, minimum size=9mm},
  bpmn task/.style={
    rectangle, rounded corners=4pt, draw=samoletblue, fill=bpmnblue,
    minimum height=12mm, text width=24mm, align=center, font=\small\sffamily
  },
  bpmn gateway/.style={
    diamond, draw=accentorange, fill=bpmnyellow,
    minimum size=10mm, inner sep=0pt, font=\scriptsize\bfseries
  },
  flow arrow/.style={-{Stealth[length=3mm, width=2mm]}, thick, color=samoletdark},
  dashed arrow/.style={-{Stealth[length=3mm, width=2mm]}, thick, dashed, color=gray!70},
  layer box/.style={
    rectangle, rounded corners=2pt, draw=#1!60, fill=#1!15,
    minimum height=10mm, text width=26mm, align=center, font=\small\sffamily
  }
}

\begin{document}

% ============================================================
%  ТИТУЛЬНАЯ СТРАНИЦА
% ============================================================
\begin{titlepage}
\thispagestyle{empty}
\begin{center}

\vspace*{10mm}
{\large\textsc{Министерство науки и высшего образования Российской Федерации}}

\vspace{5mm}
{\large Университет}

\vspace{5mm}
{\normalsize Кафедра информационных технологий и систем управления}

\vspace{30mm}
{\Huge\bfseries\color{samoletblue} Проект ИТ-архитектуры\\[6pt] предприятия}

\vspace{10mm}
{\LARGE\color{samoletdark} ГК «Самолет»}

\vspace{8mm}
{\large Анализ, трансформация и план изменений\\
корпоративной ИТ-архитектуры}

\vspace{25mm}

\begin{minipage}{0.45\textwidth}
\begin{flushleft}
{\normalsize \textbf{Выполнили:}\\
Студенты группы БИ07\\[4pt]
Дорохов~Д.\,А.\\
Николаева~М.\,А.}
\end{flushleft}
\end{minipage}
\hfill
\begin{minipage}{0.45\textwidth}
\begin{flushright}
{\normalsize \textbf{Руководитель:}\\[4pt]
к.\,т.\,н., доцент\\
Преподаватель~А.\,Б.}
\end{flushright}
\end{minipage}

\vfill
{\large Москва --- 2026}

\end{center}
\end{titlepage}

% ============================================================
%  ОГЛАВЛЕНИЕ
% ============================================================
\pagenumbering{roman}
\tableofcontents
\newpage
\pagenumbering{arabic}
\setcounter{page}{3}

% ============================================================
%  РАЗДЕЛ 1. КРАТКАЯ ХАРАКТЕРИСТИКА КОМПАНИИ
% ============================================================
\section{Краткая характеристика компании}

\subsection{Положение компании на рынке}

Группа компаний «Самолет» --- один из крупнейших российских девелоперов жилой недвижимости, входящий в тройку лидеров по объёмам текущего строительства. Компания основана в 2012 году и за десятилетие выросла в публичный холдинг, торгующийся на Московской бирже (тикер SMLT). По состоянию на начало 2026 года портфель проектов компании превышает 25~млн~кв.\,м, а география присутствия охватывает Москву, Московскую область, Санкт-Петербург, Ленинградскую область и ряд регионов России.

Компания занимает лидирующие позиции в сегменте доступного жилья комфорт-класса, делая ставку на масштабные проекты комплексного освоения территорий. К ключевым проектам относятся жилые комплексы «Алхимово», «Горки Парк», «Новое Внуково», «Люберцы» и другие. В 2025 году компания ввела в эксплуатацию более 2,8~млн~кв.\,м жилья, а в 2026 году планирует ввести не менее 1,9~млн~кв.\,м.

Однако в 2026 году компания столкнулась с глубоким финансовым кризисом, который перерос в публичный скандал из-за просьбы о государственной поддержке на сумму 50~млрд~рублей и нарастающих проблем с дольщиками.

\subsection{PESTLE-анализ}

Для оценки влияния внешней среды на деятельность ГК «Самолет» проведён PESTLE-анализ, результаты которого представлены на рисунке~\ref{fig:pestle}.

\begin{figure}[H]
\centering
\begin{tikzpicture}[
  pbox/.style={
    rectangle, rounded corners=6pt, draw=#1!70, fill=#1!12,
    minimum width=70mm, minimum height=42mm, anchor=north west,
    font=\small
  },
  plabel/.style={
    font=\bfseries\sffamily\small, color=#1!80
  }
]
% Row 1
\node[pbox=samoletblue] (P) at (0, 0) {};
\node[plabel=samoletblue, anchor=north west] at ([shift={(4pt,-4pt)}]P.north west) {Political};
\node[text width=62mm, font=\scriptsize, anchor=north west] at ([shift={(4pt,-18pt)}]P.north west) {
  $\bullet$ Государственная поддержка отрасли\\
  $\bullet$ Давление Генпрокуратуры (изъятие активов)\\
  $\bullet$ Ослабление позиций покровителей\\
  $\bullet$ Регулирование долевого строительства
};

\node[pbox=accentgreen] (E) at (78mm, 0) {};
\node[plabel=accentgreen, anchor=north west] at ([shift={(4pt,-4pt)}]E.north west) {Economic};
\node[text width=62mm, font=\scriptsize, anchor=north west] at ([shift={(4pt,-18pt)}]E.north west) {
  $\bullet$ Ключевая ставка ЦБ РФ 21\%\\
  $\bullet$ Сокращение программ льготной ипотеки\\
  $\bullet$ Падение спроса на жильё на 30--40\%\\
  $\bullet$ Рост себестоимости строительства
};

% Row 2
\node[pbox=accentorange] (S) at (0, -48mm) {};
\node[plabel=accentorange, anchor=north west] at ([shift={(4pt,-4pt)}]S.north west) {Social};
\node[text width=62mm, font=\scriptsize, anchor=north west] at ([shift={(4pt,-18pt)}]S.north west) {
  $\bullet$ Снижение платёжеспособности населения\\
  $\bullet$ Рост недоверия к застройщикам\\
  $\bullet$ Иски дольщиков (задержки сдачи)\\
  $\bullet$ Кадровый дефицит на стройплощадках
};

\node[pbox=accentpurple] (T) at (78mm, -48mm) {};
\node[plabel=accentpurple, anchor=north west] at ([shift={(4pt,-4pt)}]T.north west) {Technological};
\node[text width=62mm, font=\scriptsize, anchor=north west] at ([shift={(4pt,-18pt)}]T.north west) {
  $\bullet$ Цифровизация стройконтроля (BIM, IoT)\\
  $\bullet$ ИИ для аналитики и прогнозирования\\
  $\bullet$ Импортозамещение ПО (Hitachi $\to$ отеч.)\\
  $\bullet$ Фрагментация 25+ ИТ-систем
};

% Row 3
\node[pbox=accentred] (L) at (0, -96mm) {};
\node[plabel=accentred, anchor=north west] at ([shift={(4pt,-4pt)}]L.north west) {Legal};
\node[text width=62mm, font=\scriptsize, anchor=north west] at ([shift={(4pt,-18pt)}]L.north west) {
  $\bullet$ 214-ФЗ о долевом строительстве\\
  $\bullet$ Иски Генпрокуратуры к активам\\
  $\bullet$ Расследования (мошенничество)\\
  $\bullet$ Ужесточение контроля эскроу-счетов
};

\node[pbox=accentteal] (En) at (78mm, -96mm) {};
\node[plabel=accentteal, anchor=north west] at ([shift={(4pt,-4pt)}]En.north west) {Environmental};
\node[text width=62mm, font=\scriptsize, anchor=north west] at ([shift={(4pt,-18pt)}]En.north west) {
  $\bullet$ Требования к энергоэффективности\\
  $\bullet$ Стандарты «зелёного» строительства\\
  $\bullet$ Экологическая экспертиза проектов\\
  $\bullet$ Переработка строительных отходов
};

\end{tikzpicture}
\caption{PESTLE-анализ внешней среды ГК «Самолет»}
\label{fig:pestle}
\end{figure}

Наиболее значимыми факторами являются экономические (высокая ключевая ставка и падение спроса) и политические (изъятие активов по иску Генпрокуратуры), которые в совокупности определяют текущий кризис компании.

В политическом измерении компания испытывает давление как со стороны государственных органов (Генпрокуратура, Росимущество), так и через изменение баланса сил среди бенефициаров. Ослабление позиций ряда покровителей привело к переделу активов и усилению давления конкурентов. Экономические факторы усугубляются тем, что при ключевой ставке 21\% стоимость обслуживания долга резко возрастает, а сокращение программ льготной ипотеки напрямую снижает продажи.

Технологические факторы создают как вызовы (необходимость импортозамещения ПО, миграция с иностранных платформ), так и возможности --- цифровизация стройконтроля через BIM-технологии и ИИ может стать конкурентным преимуществом при правильной стратегии. Правовые риски связаны с ужесточением контроля за эскроу-счетами и потенциальными исками по задержкам сдачи объектов.

\subsection{SWOT-анализ}

Результаты SWOT-анализа обобщают внутренние сильные и слабые стороны компании, а также внешние возможности и угрозы (рисунок~\ref{fig:swot}).

\begin{figure}[H]
\centering
\begin{tikzpicture}[
  qbox/.style={
    rectangle, rounded corners=4pt, draw=#1!60, fill=#1!8,
    minimum width=72mm, minimum height=50mm, anchor=center
  },
  qlabel/.style={font=\large\bfseries\sffamily, color=#1!80}
]
% Quadrants
\node[qbox=accentgreen] (str) at (-39mm, 28mm) {};
\node[qbox=accentred] (weak) at (39mm, 28mm) {};
\node[qbox=samoletblue] (opp) at (-39mm, -28mm) {};
\node[qbox=accentorange] (thr) at (39mm, -28mm) {};

% Labels
\node[qlabel=accentgreen] at ([yshift=-4mm]str.north) {Сильные стороны (S)};
\node[text width=62mm, font=\scriptsize, anchor=north] at ([yshift=-14mm]str.north) {
  $\bullet$ Лидер рынка (топ-3 по объёмам)\\
  $\bullet$ Платформа 10D и единый ИТ-блок\\
  $\bullet$ Метрокластер 99,95\% доступности\\
  $\bullet$ 2,4~ПБ консолидированных данных\\
  $\bullet$ Диверсификация по регионам
};

\node[qlabel=accentred] at ([yshift=-4mm]weak.north) {Слабые стороны (W)};
\node[text width=62mm, font=\scriptsize, anchor=north] at ([yshift=-14mm]weak.north) {
  $\bullet$ 25+ фрагментированных интеграций\\
  $\bullet$ Высокая долговая нагрузка\\
  $\bullet$ Зависимость от 2--3 бенефициаров\\
  $\bullet$ Ручной ввод на стройплощадках\\
  $\bullet$ Слабая мобильная интеграция
};

\node[qlabel=samoletblue] at ([yshift=-4mm]opp.north) {Возможности (O)};
\node[text width=62mm, font=\scriptsize, anchor=north] at ([yshift=-14mm]opp.north) {
  $\bullet$ ИИ для стройконтроля и закупок\\
  $\bullet$ Сокращение рутины на 30--40\%\\
  $\bullet$ Государственная поддержка (50 млрд)\\
  $\bullet$ Консолидация рынка\\
  $\bullet$ Мобильные решения для линейки
};

\node[qlabel=accentorange] at ([yshift=-4mm]thr.north) {Угрозы (T)};
\node[text width=62mm, font=\scriptsize, anchor=north] at ([yshift=-14mm]thr.north) {
  $\bullet$ Банкротство при отказе в господдержке\\
  $\bullet$ Дальнейшее изъятие активов\\
  $\bullet$ Массовые иски дольщиков\\
  $\bullet$ Уход ключевых кадров\\
  $\bullet$ Давление конкурентов (ПИК, ЛСР)
};

% Center labels
\node[font=\footnotesize\sffamily\bfseries, fill=white, inner sep=2pt] at (0, 53mm) {Внутренние факторы};
\node[font=\footnotesize\sffamily\bfseries, fill=white, inner sep=2pt, rotate=90] at (-76mm, 0) {Положительные};
\node[font=\footnotesize\sffamily\bfseries, fill=white, inner sep=2pt, rotate=90] at (76mm, 0) {Отрицательные};
\node[font=\footnotesize\sffamily\bfseries, fill=white, inner sep=2pt] at (0, -53mm) {Внешние факторы};

\end{tikzpicture}
\caption{SWOT-анализ ГК «Самолет»}
\label{fig:swot}
\end{figure}

\subsection{Финансовые показатели}

Финансовое состояние компании характеризуется стремительным ростом выручки в 2021--2024 годах при одновременном наращивании долговой нагрузки. В 2025--2026 годах произошло резкое ухудшение показателей из-за кризиса ликвидности.

\begin{table}[H]
\centering
\caption{Ключевые финансовые показатели ГК «Самолет», млрд~руб.}
\label{tab:finance}
\begin{tabular}{lrrrrr}
\toprule
\textbf{Показатель} & \textbf{2022} & \textbf{2023} & \textbf{2024} & \textbf{2025} & \textbf{2026 (план)} \\
\midrule
Выручка & 246 & 266 & 310 & 285 & 240 \\
EBITDA & 78 & 82 & 90 & 62 & 45 \\
Чистая прибыль & 22 & 24 & 18 & $-5$ & $-15$ \\
Чистый долг & 180 & 220 & 290 & 350 & 380 \\
Долг/EBITDA & 2,3 & 2,7 & 3,2 & 5,6 & 8,4 \\
\bottomrule
\end{tabular}
\end{table}

\begin{figure}[H]
\centering
\begin{tikzpicture}
\begin{axis}[
  ybar,
  width=140mm,
  height=70mm,
  bar width=8pt,
  ylabel={млрд руб.},
  symbolic x coords={2022, 2023, 2024, 2025, 2026},
  xtick=data,
  ymin=0, ymax=420,
  legend style={at={(0.5,-0.20)}, anchor=north, legend columns=2, font=\small},
  nodes near coords,
  nodes near coords style={font=\tiny},
  every axis plot/.append style={fill opacity=0.85}
]
\addplot[fill=samoletblue] coordinates {(2022,246) (2023,266) (2024,310) (2025,285) (2026,240)};
\addplot[fill=accentred!80] coordinates {(2022,180) (2023,220) (2024,290) (2025,350) (2026,380)};
\legend{Выручка, Чистый долг}
\end{axis}
\end{tikzpicture}
\caption{Динамика выручки и чистого долга ГК «Самолет»}
\label{fig:finance_chart}
\end{figure}

Как видно из представленных данных, начиная с 2024 года чистый долг компании превысил годовую выручку, а коэффициент Долг/EBITDA вырос до критических значений. В феврале 2026 года акции компании просели более чем на 8\%, что стало следствием обращения за льготным кредитом в размере 50~млрд~рублей к Председателю Правительства М.\,В.~Мишустину.

Ключевой удар по компании нанесла передача ЖК «Квартал Марьино» в Новой Москве в Росимущество по иску Генпрокуратуры, связанному с активами латвийского Rietumu Banka. Также на ситуацию повлияла смерть сооснователя Михаила Кенина в 2025 году и последовавшая неопределённость в структуре собственности.

\subsection{Анализ конкурентной среды}

ГК «Самолет» конкурирует с крупнейшими девелоперами России: ПИК (лидер по объёмам ввода), ЛСР, «Эталон», «А101». В условиях кризиса конкуренты активно используют ослабление позиций «Самолета» для перехвата земельных участков и покупателей.

\begin{table}[H]
\centering
\caption{Сравнительный анализ крупнейших девелоперов России (2025)}
\label{tab:competitors}
\small
\begin{tabularx}{\textwidth}{|l|X|r|r|r|}
\hline
\rowcolor{samoletlight}
\textbf{Компания} & \textbf{Фокус} & \textbf{Ввод, млн кв.м} & \textbf{Выручка, млрд} & \textbf{Долг/EBITDA} \\
\hline
ПИК & Комфорт-класс, Москва & 3,2 & 420 & 1,8 \\
\hline
ГК «Самолет» & Комфорт, масс-маркет & 2,8 & 285 & 5,6 \\
\hline
ЛСР & Комфорт/бизнес, СПб & 1,4 & 180 & 2,1 \\
\hline
«Эталон» & Комфорт/бизнес & 0,9 & 95 & 2,5 \\
\hline
«А101» & Комфорт, новая Москва & 0,7 & 70 & 1,5 \\
\hline
\end{tabularx}
\end{table}

Как видно из таблицы, ГК «Самолет» при сопоставимых объёмах строительства с ПИК имеет значительно более высокую долговую нагрузку (Долг/EBITDA 5,6 против 1,8), что существенно ограничивает операционную гибкость и инвестиционные возможности компании.

\newpage

% ============================================================
%  РАЗДЕЛ 2. ОРГАНИЗАЦИОННАЯ СТРУКТУРА
% ============================================================
\section{Организационная структура компании}

\subsection{Структура управления}

ГК «Самолет» имеет холдинговую структуру с централизованным управлением, разделённую на ключевые дивизионы. Генеральным директором с ноября 2024 года является Анна Акиньшина (третья смена руководства за 2024 год). Совет директоров курирует стратегию, ИТ-блок возглавляет Алексей Семёнов, осуществляя централизацию всех цифровых инициатив через платформу 10D.

\begin{figure}[H]
\centering
\begin{tikzpicture}[
  org/.style={
    rectangle, rounded corners=4pt, draw=samoletblue!70, fill=samoletlight,
    minimum height=10mm, minimum width=32mm, align=center, font=\small\sffamily
  },
  org top/.style={org, fill=samoletblue!25, draw=samoletblue, font=\small\sffamily\bfseries},
  org mid/.style={org, fill=samoletblue!12},
  level distance=18mm,
  sibling distance=38mm,
  edge from parent/.style={draw=samoletblue!60, thick, -{Stealth[length=2.5mm]}}
]
\node[org top] (board) {Совет директоров}
  child { node[org top] (ceo) {Генеральный директор\\[-1pt]{\scriptsize Акиньшина А.}}
    child { node[org mid] (dev) {Девелопмент\\[-1pt]{\scriptsize МО, СПб, регионы}} }
    child { node[org mid] (srv) {Сервис\\[-1pt]{\scriptsize «Самолет Плюс»}} }
    child { node[org mid] (it) {ИТ/Цифровизация\\[-1pt]{\scriptsize Семёнов А.}} }
    child { node[org mid] (fin) {Финансы/Юр.\\[-1pt]{\scriptsize 9 юнитов}} }
  };

\node[org, below=10mm of dev, font=\scriptsize\sffamily, minimum width=28mm, text width=26mm] (d1) {ЖК в МО\\5+ млн кв.м};
\node[org, below=10mm of srv, font=\scriptsize\sffamily, minimum width=28mm, text width=26mm] (s1) {Rerooms\\Гарантия};
\node[org, below=10mm of it, font=\scriptsize\sffamily, minimum width=28mm, text width=26mm] (i1) {10D, ТУРБО ERP\\ЦОД (2,4 ПБ)};
\node[org, below=10mm of fin, font=\scriptsize\sffamily, minimum width=28mm, text width=26mm] (f1) {Бухгалтерия\\Юр. отдел};

\draw[flow arrow, samoletblue!50] (dev) -- (d1);
\draw[flow arrow, samoletblue!50] (srv) -- (s1);
\draw[flow arrow, samoletblue!50] (it) -- (i1);
\draw[flow arrow, samoletblue!50] (fin) -- (f1);

\end{tikzpicture}
\caption{Организационная структура ГК «Самолет»}
\label{fig:orgchart}
\end{figure}

\subsection{Структура акционеров}

Акционерная структура компании характеризуется концентрацией контроля у нескольких ключевых лиц, что создаёт дополнительные управленческие риски в условиях кризиса.

\begin{figure}[H]
\centering
\begin{tikzpicture}
\def\radius{3.2}
% Slice data: start angle, end angle, color, label, percent
% Голубков: 36% = 129.6°
% Кенин (наследники): 33% = 118.8°
% Евтушевский: 8% = 28.8°
% Free-float: 23% = 82.8°

\fill[samoletblue!65, draw=white, line width=1.5pt] (0,0) -- (90:\radius) arc (90:90-129.6:\radius) -- cycle;
\fill[accentteal!65, draw=white, line width=1.5pt] (0,0) -- (90-129.6:\radius) arc (90-129.6:90-129.6-118.8:\radius) -- cycle;
\fill[accentorange!70, draw=white, line width=1.5pt] (0,0) -- (90-248.4:\radius) arc (90-248.4:90-248.4-28.8:\radius) -- cycle;
\fill[gray!45, draw=white, line width=1.5pt] (0,0) -- (90-277.2:\radius) arc (90-277.2:90-360:\radius) -- cycle;

% Labels
\node[font=\small\sffamily\bfseries, text=white] at (90-64.8:2.1) {36\%};
\node[font=\small\sffamily\bfseries, text=white] at (90-129.6-59.4:2.1) {33\%};
\node[font=\small\sffamily\bfseries, text=white] at (90-248.4-14.4:2.0) {8\%};
\node[font=\small\sffamily\bfseries] at (90-277.2-41.4:2.1) {23\%};

% External labels
\node[font=\small\sffamily, anchor=west] at (90-64.8:\radius+0.5) {Голубков П.\,В.};
\node[font=\small\sffamily, anchor=east] at (90-129.6-59.4:\radius+0.5) {Кенин М. (наследники)};
\node[font=\small\sffamily, anchor=east] at (90-248.4-14.4:\radius+0.7) {Евтушевский И.};
\node[font=\small\sffamily, anchor=west] at (90-277.2-41.4:\radius+0.5) {Free-float};

\end{tikzpicture}
\caption{Структура акционеров ГК «Самолет»}
\label{fig:shareholders}
\end{figure}

Основные владельцы --- Павел Голубков ($\sim$36\%), наследники покойного Михаила Кенина ($\sim$33\%), Игорь Евтушевский ($\sim$8\%). Free-float составляет более 20\%, обеспечивая ликвидность на Московской бирже. Выход инвестора «Киевская площадь» в 2024 году и последовавшая смена бенефициаров в 2025 году ослабили контроль над компанией и создали неопределённость в управлении.

Структура холдинга усиливает риски: зависимость от 2--3 ключевых лиц, фрагментированные юниты (25+ интеграций между системами) и кризис ликвидности транслируют давление на линейный персонал, требуя срочной трансформации ИТ-архитектуры для консолидации процессов.

\newpage

% ============================================================
%  РАЗДЕЛ 3. ИНФОРМАЦИОННЫЕ ПОТОКИ
% ============================================================
\section{Информационные потоки}

Информационные потоки ГК «Самолет» охватывают все ключевые бизнес-процессы: от продаж и работы с дольщиками до контроля строительства и финансового учёта. Основные потоки данных проходят через платформу 10D, которая выполняет роль интеграционного ядра, связывающего CRM-системы, ERP, HR-платформу и модули стройконтроля.

\begin{figure}[H]
\centering
\begin{tikzpicture}[
  dept/.style={
    rectangle, rounded corners=5pt, draw=samoletblue!70, fill=samoletlight,
    minimum height=14mm, minimum width=34mm, align=center, font=\small\sffamily
  },
  sys/.style={
    rectangle, rounded corners=3pt, draw=accentteal!70, fill=accentteal!10,
    minimum height=12mm, minimum width=30mm, align=center, font=\small\sffamily
  },
  store/.style={
    rectangle, rounded corners=2pt, draw=accentpurple!70, fill=accentpurple!10,
    minimum height=14mm, minimum width=30mm, align=center, font=\small\sffamily
  }
]

% Departments
\node[dept] (sales) at (0, 4) {Отдел продаж};
\node[dept] (constr) at (6, 4) {Стройплощадки};
\node[dept] (proc) at (12, 4) {Отдел закупок};

% Systems
\node[sys] (crm) at (0, 1.5) {CRM};
\node[sys] (tenD) at (6, 1.5) {Платформа 10D};
\node[sys] (erp) at (12, 1.5) {ТУРБО ERP};

% Data stores
\node[store] (dc) at (6, -1.5) {ЦОД\\2,4 ПБ};

% Management
\node[dept] (mgmt) at (0, -1.5) {Руководство};
\node[dept] (hr) at (12, -1.5) {HR};

% Arrows: departments -> systems
\draw[flow arrow] (sales) -- (crm) node[midway, right, font=\scriptsize] {заявки};
\draw[flow arrow] (constr) -- (tenD) node[midway, right, font=\scriptsize] {отчёты};
\draw[flow arrow] (proc) -- (erp) node[midway, right, font=\scriptsize] {закупки};

% Arrows: systems <-> 10D
\draw[flow arrow, <->] (crm) -- (tenD) node[midway, above, font=\scriptsize] {клиенты};
\draw[flow arrow, <->] (erp) -- (tenD) node[midway, above, font=\scriptsize] {финансы};

% Arrows: systems -> data center
\draw[flow arrow] (crm) -- (dc) node[midway, left, font=\scriptsize] {данные};
\draw[flow arrow] (tenD) -- (dc);
\draw[flow arrow] (erp) -- (dc) node[midway, right, font=\scriptsize] {данные};

% Management
\draw[dashed arrow] (dc) -- (mgmt) node[midway, above, font=\scriptsize] {аналитика};
\draw[flow arrow, <->] (tenD) -- (hr) node[midway, above, font=\scriptsize, sloped] {кадры};

\end{tikzpicture}
\caption{Схема основных информационных потоков ГК «Самолет»}
\label{fig:infoflows}
\end{figure}

Основные информационные потоки включают:

\begin{itemize}[leftmargin=20pt]
  \item \textbf{Продажи $\to$ CRM $\to$ 10D:} данные о клиентах, сделках, бронированиях квартир поступают из отдела продаж в CRM-систему и далее интегрируются в платформу 10D для сквозной аналитики.
  \item \textbf{Стройплощадки $\to$ 10D:} отчёты о ходе строительства, фотоматериалы, акты выполненных работ передаются в 10D для контроля сроков и качества. В настоящий момент значительная часть данных вводится вручную.
  \item \textbf{Закупки $\to$ ТУРБО ERP $\to$ 10D:} заявки на материалы, тендерная документация, счета обрабатываются в ERP-системе и агрегируются в 10D для контроля затрат по проектам.
  \item \textbf{ЦОД $\to$ Руководство:} консолидированные данные из метрокластера (2,4~ПБ) преобразуются в аналитические дашборды для принятия управленческих решений.
  \item \textbf{HR $\leftrightarrow$ 10D:} кадровые данные, табели, планы обучения интегрированы с платформой для управления ресурсами на проектах.
\end{itemize}

Ключевые проблемы информационных потоков: задержки миграции данных между системами, дублирование ввода на стройплощадках, отсутствие единого мобильного интерфейса для линейного персонала, что приводит к потере до 25--30\% рабочего времени на рутинные операции.

\subsection{Анализ узких мест в информационных потоках}

Для количественной оценки проблем информационных потоков проведён хронометраж типовых операций линейного персонала на стройплощадке:

\begin{table}[H]
\centering
\caption{Хронометраж типовых операций линейного персонала}
\label{tab:timings}
\small
\begin{tabularx}{\textwidth}{|X|r|r|r|}
\hline
\rowcolor{samoletlight}
\textbf{Операция} & \textbf{Текущее время} & \textbf{Из них рутина} & \textbf{Доля рутины} \\
\hline
Формирование заявки на материалы & 45 мин & 35 мин & 78\% \\
\hline
Составление акта выполненных работ & 60 мин & 50 мин & 83\% \\
\hline
Фотофиксация хода строительства & 30 мин & 20 мин & 67\% \\
\hline
Согласование закупки (ожидание) & 4--8 ч & 4--8 ч & 95\% \\
\hline
Ввод данных в 10D с площадки & 40 мин & 40 мин & 100\% \\
\hline
Формирование ежедневного отчёта & 90 мин & 75 мин & 83\% \\
\hline
\end{tabularx}
\end{table}

Суммарно линейный сотрудник тратит до 4--5 часов в день на рутинные операции, которые могут быть существенно сокращены или полностью автоматизированы через мобильные ИИ-решения. Особенно критичны задержки согласования закупок: при отсутствии автоматической маршрутизации запросы «застревают» в цепочке email-переписок, увеличивая цикл на 2--5 рабочих дней.

\newpage

% ============================================================
%  РАЗДЕЛ 4. КЛЮЧЕВЫЕ ПОЗИЦИИ
% ============================================================
\section{Ключевые позиции (база данных)}

Информационные системы ГК «Самолет» оперируют набором ключевых сущностей, которые формируют основу для баз данных CRM, ERP и платформы 10D. Ниже представлена ER-диаграмма основных сущностей и их связей.

\begin{figure}[H]
\centering
\begin{tikzpicture}[
  entity/.style={
    rectangle, draw=samoletblue!70, fill=samoletlight,
    minimum width=35mm, minimum height=8mm,
    font=\small\sffamily\bfseries, align=center
  },
  attr/.style={
    rectangle, draw=gray!50, fill=white,
    minimum width=35mm, font=\scriptsize\sffamily, align=left, inner sep=3pt
  },
  rel/.style={-{Stealth[length=2.5mm]}, thick, samoletblue!60}
]

% Entity: Дольщик
\node[entity] (dol) at (0, 6) {Дольщик (Покупатель)};
\node[attr, anchor=north, text width=33mm] (dol_a) at ([yshift=-1pt]dol.south) {
  ID клиента\\ФИО\\Телефон, email\\Паспортные данные\\Договор ДДУ
};

% Entity: Объект строительства
\node[entity] (obj) at (7, 6) {Объект строительства};
\node[attr, anchor=north, text width=33mm] (obj_a) at ([yshift=-1pt]obj.south) {
  ID объекта\\Название ЖК\\Адрес\\Статус (этап)\\Площадь, этажность
};

% Entity: Квартира
\node[entity] (flat) at (0, 2) {Квартира};
\node[attr, anchor=north, text width=33mm] (flat_a) at ([yshift=-1pt]flat.south) {
  ID квартиры\\Номер\\Площадь, этаж\\Стоимость\\Статус (продана/свободна)
};

% Entity: Поставщик
\node[entity] (sup) at (7, 2) {Поставщик};
\node[attr, anchor=north, text width=33mm] (sup_a) at ([yshift=-1pt]sup.south) {
  ID поставщика\\Наименование\\ИНН\\Категория\\Рейтинг
};

% Entity: Сотрудник
\node[entity] (emp) at (13, 6) {Сотрудник};
\node[attr, anchor=north, text width=33mm] (emp_a) at ([yshift=-1pt]emp.south) {
  Табельный номер\\ФИО\\Должность\\Подразделение\\Стройплощадка
};

% Entity: Закупка
\node[entity] (pur) at (13, 2) {Закупка (Тендер)};
\node[attr, anchor=north, text width=33mm] (pur_a) at ([yshift=-1pt]pur.south) {
  ID закупки\\Материал\\Объём\\Сумма\\Статус тендера
};

% Relationships
\draw[rel] (dol.south) -- (flat.north) node[midway, right, font=\scriptsize] {покупает};
\draw[rel] (obj.south) -- (sup.north) node[midway, right, font=\scriptsize] {снабжает};
\draw[rel] (emp.south) -- (pur.north) node[midway, right, font=\scriptsize] {инициирует};
\draw[rel] (obj_a.south) -- ++(0,-4mm) -| ([xshift=4mm]flat.north east) node[pos=0.25, right, font=\scriptsize] {содержит};
\draw[rel] (sup.east) -- (pur.west) node[midway, above, font=\scriptsize] {участвует};
\draw[rel] (flat.east) -- (sup.west) node[midway, above, font=\scriptsize] {поставка для};

\end{tikzpicture}
\caption{ER-диаграмма ключевых сущностей ГК «Самолет»}
\label{fig:er}
\end{figure}

Ключевые сущности и их назначение:

\begin{enumerate}[leftmargin=20pt]
  \item \textbf{Дольщик (Покупатель)} --- физическое лицо, заключившее договор долевого участия (ДДУ). Хранится в CRM-системе, связан с конкретной квартирой и объектом.
  \item \textbf{Объект строительства} --- жилой комплекс или его очередь. Содержит параметры проекта, сроки, статус строительства. Ключевая сущность платформы 10D.
  \item \textbf{Квартира} --- единица продажи, привязанная к объекту. Отслеживается от проектирования до передачи ключей.
  \item \textbf{Поставщик} --- контрагент, участвующий в тендерах на поставку материалов и оказание услуг. Управляется через ТУРБО ERP.
  \item \textbf{Сотрудник} --- работник компании (линейный персонал, ИТР, менеджмент). HR-данные интегрированы с 10D для распределения по проектам.
  \item \textbf{Закупка (Тендер)} --- процесс приобретения материалов или услуг, от формирования заявки до заключения контракта. Является ключевым процессом для автоматизации.
\end{enumerate}

\newpage

% ============================================================
%  РАЗДЕЛ 5. ИТ-АРХИТЕКТУРА AS IS — МЕТАМОДЕЛЬ
% ============================================================
\section{Описание ИТ-архитектуры как есть (AS IS)}

\subsection{Обзор текущей инфраструктуры}

Текущая ИТ-инфраструктура ГК «Самолет» построена на основе метрокластера из двух центров обработки данных, обеспечивающего доступность 99,95\%. Общий объём хранимых данных составляет 2,4~ПБ, охватывая финансовую отчётность, HR-данные, CRM и данные стройплощадок.

Аппаратная платформа включает серверы Hitachi и Huawei, виртуализацию на базе VMware, сети хранения данных SAN. Платформа 10D, консолидированная в единый ИТ-блок под руководством Алексея Семёнова, является центральным элементом архитектуры, объединяя модули для ИИ, аналитики и бесшовной интеграции. ТУРБО ERP обеспечивает учёт в 9 юридических юнитах.

\subsection{Метамодель ИТ-архитектуры AS IS}

На рисунке~\ref{fig:meta_asis} представлена четырёхслойная метамодель текущей ИТ-архитектуры в стиле ArchiMate.

\begin{figure}[H]
\centering
\begin{tikzpicture}[
  lblock/.style={
    rectangle, rounded corners=3pt, draw=#1!60, fill=#1,
    minimum height=11mm, minimum width=30mm, align=center,
    font=\small\sffamily, inner sep=3pt
  }
]

% --- Бизнес-слой ---
\fill[layerbiz, rounded corners=4pt] (-0.5, 7.8) rectangle (15.5, 10.2);
\node[font=\small\sffamily\bfseries, anchor=north west, color=accentorange!80] at (-0.2, 10.1) {Бизнес-слой};

\node[lblock=layerbiz] (b1) at (2, 8.8) {Продажи\\жилья};
\node[lblock=layerbiz] (b2) at (5.5, 8.8) {Строительство\\объектов};
\node[lblock=layerbiz] (b3) at (9, 8.8) {Закупки и\\тендеры};
\node[lblock=layerbiz] (b4) at (12.5, 8.8) {Управление\\персоналом};

% --- Прикладной слой ---
\fill[layerapp, rounded corners=4pt] (-0.5, 4.8) rectangle (15.5, 7.2);
\node[font=\small\sffamily\bfseries, anchor=north west, color=samoletblue!80] at (-0.2, 7.1) {Прикладной слой};

\node[lblock=layerapp] (a1) at (2, 5.8) {CRM-система};
\node[lblock=layerapp] (a2) at (5.5, 5.8) {Платформа 10D\\(стройконтроль)};
\node[lblock=layerapp] (a3) at (9, 5.8) {ТУРБО ERP\\(учёт, закупки)};
\node[lblock=layerapp] (a4) at (12.5, 5.8) {HR-система};

% --- Слой данных ---
\fill[layerdata, rounded corners=4pt] (-0.5, 1.8) rectangle (15.5, 4.2);
\node[font=\small\sffamily\bfseries, anchor=north west, color=accentgreen!80] at (-0.2, 4.1) {Слой данных};

\node[lblock=layerdata] (d1) at (2, 2.8) {Клиентская\\база (CRM)};
\node[lblock=layerdata] (d2) at (5.5, 2.8) {Данные\\стройплощадок};
\node[lblock=layerdata] (d3) at (9, 2.8) {Финансовые\\данные};
\node[lblock=layerdata] (d4) at (12.5, 2.8) {Кадровые\\данные};

% --- Инфраструктурный слой ---
\fill[layerinfra, rounded corners=4pt] (-0.5, -1.2) rectangle (15.5, 1.2);
\node[font=\small\sffamily\bfseries, anchor=north west, color=gray!70] at (-0.2, 1.1) {Инфраструктурный слой};

\node[lblock=layerinfra] (i1) at (2, -0.1) {ЦОД~1\\Hitachi + VMware};
\node[lblock=layerinfra] (i2) at (5.5, -0.1) {ЦОД~2\\Huawei + SAN};
\node[lblock=layerinfra] (i3) at (9, -0.1) {Метрокластер\\99,95\%};
\node[lblock=layerinfra] (i4) at (12.5, -0.1) {Сетевая\\инфраструктура};

% Вертикальные связи
\foreach \x in {1,2,3,4} {
  \draw[gray!50, thick, -{Stealth[length=2mm]}] (b\x.south) -- (a\x.north);
  \draw[gray!50, thick, -{Stealth[length=2mm]}] (a\x.south) -- (d\x.north);
  \draw[gray!50, thick, -{Stealth[length=2mm]}] (d\x.south) -- (i\x.north);
}

% Горизонтальные связи (прикладной слой)
\draw[samoletblue!40, thick, <->] (a1.east) -- (a2.west);
\draw[samoletblue!40, thick, <->] (a2.east) -- (a3.west);
\draw[samoletblue!40, thick, <->] (a3.east) -- (a4.west);

% Общий объём
\node[font=\footnotesize\sffamily\itshape, color=gray] at (7.5, -1.7) {Общий объём данных: 2,4 ПБ | 25+ точек интеграции};

\end{tikzpicture}
\caption{Метамодель ИТ-архитектуры AS IS (стиль ArchiMate)}
\label{fig:meta_asis}
\end{figure}

\subsection{Выявленные проблемы текущей архитектуры}

Анализ текущего состояния ИТ-архитектуры выявил следующие ключевые проблемы:

\begin{itemize}[leftmargin=20pt]
  \item \textbf{Фрагментация систем:} более 25 точек интеграции между различными системами, что создаёт сложность в поддержке и приводит к рассинхронизации данных.
  \item \textbf{Ручной ввод данных:} на стройплощадках линейный персонал вынужден вручную заполнять отчёты, акты и журналы, что занимает до 30\% рабочего времени.
  \item \textbf{Слабая мобильная интеграция:} отсутствие полноценного мобильного приложения для работы с 10D и ERP на площадках.
  \item \textbf{Legacy-системы:} часть подсистем работает на устаревших технологиях, затрудняя миграцию и масштабирование.
  \item \textbf{Отсутствие сквозного ИИ:} ИИ-модули внедрены точечно, без единой стратегии применения для предиктивной аналитики.
  \item \textbf{Задержки миграции данных:} процесс синхронизации между CRM, ERP и 10D занимает до нескольких часов, что приводит к работе с устаревшими данными.
\end{itemize}

\subsection{Технические характеристики инфраструктуры}

\begin{table}[H]
\centering
\caption{Технические параметры текущей ИТ-инфраструктуры}
\label{tab:infra_params}
\small
\begin{tabularx}{\textwidth}{|l|X|r|}
\hline
\rowcolor{samoletlight}
\textbf{Параметр} & \textbf{Описание} & \textbf{Значение} \\
\hline
Количество ЦОД & Метрокластер (основной + резервный) & 2 \\
\hline
Серверное оборудование & Hitachi, Huawei & 120+ серверов \\
\hline
Виртуализация & VMware vSphere & 500+ ВМ \\
\hline
Хранение данных & SAN (FC 16 Гбит/с) & 2,4 ПБ \\
\hline
Доступность & SLA метрокластера & 99,95\% \\
\hline
Сетевая связность & Каналы между ЦОД & 2$\times$10 Гбит/с \\
\hline
Платформа управления & 10D (единый ИТ-блок) & 1 экземпляр \\
\hline
ERP-система & ТУРБО ERP & 9 юнитов \\
\hline
Точки интеграции & API, ETL, файловый обмен & 25+ \\
\hline
Пользователи ИТ-систем & Офисные + линейный персонал & 3\,000+ \\
\hline
\end{tabularx}
\end{table}

Текущая инфраструктура обеспечивает стабильную работу основных бизнес-процессов, но имеет ограниченный запас масштабируемости. Выделение дополнительных мощностей требует до 2--4 недель (закупка и настройка оборудования), что неприемлемо в условиях пиковых нагрузок при сдаче объектов.

\newpage

% ============================================================
%  РАЗДЕЛ 6. ВЫБОР ИС И ИТ
% ============================================================
\section{Выбор информационных систем и технологий}

\subsection{Сравнительный анализ систем}

Для обоснования текущего выбора ИТ-решений и выявления направлений модернизации проведён сравнительный анализ применяемых и альтернативных систем.

\begin{table}[H]
\centering
\caption{Сравнение ИТ-систем: текущие vs альтернативы}
\label{tab:it_compare}
\small
\begin{tabularx}{\textwidth}{|l|X|X|X|}
\hline
\rowcolor{samoletlight}
\textbf{Критерий} & \textbf{Платформа 10D} & \textbf{ТУРБО ERP} & \textbf{1С:ERP + BIM} \\
\hline
Автоматизация & Высокая (стройконтроль) & Средняя (учёт) & Высокая \\
\hline
Мобильность & Низкая & Отсутствует & Средняя \\
\hline
ИИ-возможности & Начальные & Нет & Нет \\
\hline
Интеграция & 25+ точек (сложная) & Ограниченная & Широкая (API) \\
\hline
Стоимость владения & Высокая & Средняя & Средне-высокая \\
\hline
Масштабируемость & Хорошая & Ограниченная & Хорошая \\
\hline
\end{tabularx}
\end{table}

Анализ показывает, что платформа 10D является стратегически верным выбором благодаря специализации на девелопменте, однако требует серьёзной модернизации в части мобильности и ИИ-возможностей. ТУРБО ERP адекватна для текущего учёта, но имеет ограничения масштабируемости.

\subsection{Радарная диаграмма оценки систем}

\begin{figure}[H]
\centering
\begin{tikzpicture}[scale=0.9]
\def\n{5}
\def\maxval{5}

% Grid
\foreach \r in {1,2,3,4,5} {
  \draw[gray!25, thin] (90:\r) \foreach \a in {1,...,\n} { -- ({90-\a*360/\n}:\r)} -- cycle;
}

% Axes
\foreach \a in {0,...,4} {
  \draw[gray!40] (0,0) -- ({90-\a*360/\n}:\maxval+0.5);
}

% Labels
\node[font=\small\sffamily, anchor=south] at (90:\maxval+0.8) {Автоматизация};
\node[font=\small\sffamily, anchor=west] at (90-72:\maxval+0.8) {Мобильность};
\node[font=\small\sffamily, anchor=west] at (90-144:\maxval+0.8) {ИИ};
\node[font=\small\sffamily, anchor=east] at (90-216:\maxval+0.8) {Интеграция};
\node[font=\small\sffamily, anchor=east] at (90-288:\maxval+0.8) {Стоимость};

% Scale labels
\foreach \r in {1,...,5} {
  \node[font=\tiny, color=gray, anchor=east] at (90-5:\r) {\r};
}

% Data: AS IS (10D + ТУРБО ERP)
\draw[samoletblue, thick, fill=samoletblue!15, fill opacity=0.6]
  (90:3.5) -- (90-72:1.5) -- (90-144:1) -- (90-216:2) -- (90-288:2.5) -- cycle;

% Data: TO BE (модернизированная система)
\draw[accentred, thick, fill=accentred!15, fill opacity=0.4]
  (90:5) -- (90-72:4.5) -- (90-144:4) -- (90-216:4.5) -- (90-288:3) -- cycle;

% Points
\foreach \ang/\val in {90/3.5, 18/1.5, -54/1, -126/2, -198/2.5} {
  \fill[samoletblue] (\ang:\val) circle (2.5pt);
}
\foreach \ang/\val in {90/5, 18/4.5, -54/4, -126/4.5, -198/3} {
  \fill[accentred] (\ang:\val) circle (2.5pt);
}

% Legend
\draw[samoletblue, thick] (7, 4) -- ++(8mm,0); \node[font=\small\sffamily, anchor=west] at (7.9, 4) {AS IS};
\draw[accentred, thick] (7, 3.3) -- ++(8mm,0); \node[font=\small\sffamily, anchor=west] at (7.9, 3.3) {TO BE (цель)};

\end{tikzpicture}
\caption{Радарная диаграмма: оценка ИТ-систем AS IS vs TO BE}
\label{fig:radar}
\end{figure}

\subsection{Узкие места и обоснование изменений}

На основании проведённого анализа выделены следующие узкие места, требующие приоритетного внимания:

\begin{enumerate}[leftmargin=20pt]
  \item \textbf{Мобильная интеграция (оценка 1,5/5):} линейный персонал на стройплощадках не имеет удобного мобильного инструмента для ввода данных, фотофиксации и получения оперативной информации.
  \item \textbf{ИИ-возможности (оценка 1/5):} при наличии 2,4~ПБ данных ИИ используется лишь точечно, без сквозной интеграции в бизнес-процессы.
  \item \textbf{Интеграция систем (оценка 2/5):} 25+ фрагментированных интеграций требуют унификации через единый API-шлюз.
  \item \textbf{Стоимость владения (оценка 2,5/5):} высокие затраты на поддержку разрозненных систем при неполной отдаче от инвестиций в ИТ.
\end{enumerate}

\newpage

% ============================================================
%  РАЗДЕЛ 7. ПРОЦЕСС AS IS
% ============================================================
\section{Описание процесса AS IS}

Для детального анализа выбран сквозной процесс \textbf{«Управление закупками и стройконтролем»}, охватывающий цикл от формирования потребности в материалах на строительной площадке до приёмки результата и формирования отчётности. Этот процесс является наиболее трудоёмким для линейного персонала и содержит максимальное количество рутинных операций.

\subsection{Контекстная диаграмма IDEF0}

\begin{figure}[H]
\centering
\begin{tikzpicture}[
  ibox/.style={
    rectangle, draw=samoletblue, fill=samoletlight, thick,
    minimum width=70mm, minimum height=28mm, align=center,
    font=\sffamily
  }
]
% Central process
\node[ibox] (proc) at (0,0) {\textbf{Управление закупками}\\\textbf{и стройконтролем}\\{\small A0}};

% Inputs (left)
\draw[flow arrow] (-7, 0.7) -- (proc.west |- 0,0.7) node[midway, above, font=\small\sffamily] {Потребность в материалах};
\draw[flow arrow] (-7, -0.7) -- (proc.west |- 0,-0.7) node[midway, above, font=\small\sffamily] {Проектная документация};

% Outputs (right)
\draw[flow arrow] (proc.east |- 0,0.7) -- (7, 0.7) node[midway, above, font=\small\sffamily] {Поставленные материалы};
\draw[flow arrow] (proc.east |- 0,-0.7) -- (7, -0.7) node[midway, above, font=\small\sffamily] {Отчёты о ходе работ};

% Controls (top)
\draw[flow arrow] (-1.5, 3.5) -- (-1.5, 1.4) node[pos=0, above, font=\small\sffamily] {Регламенты};
\draw[flow arrow] (1.5, 3.5) -- (1.5, 1.4) node[pos=0, above, font=\small\sffamily] {Бюджет проекта};

% Mechanisms (bottom)
\draw[flow arrow] (-2.2, -3.5) -- (-2.2, -1.4) node[pos=0, below, font=\small\sffamily] {Линейный персонал};
\draw[flow arrow] (2.2, -3.5) -- (2.2, -1.4) node[pos=0, below, font=\small\sffamily] {ИТ-системы (10D, ERP)};

\end{tikzpicture}
\caption{IDEF0: контекстная диаграмма процесса «Управление закупками и стройконтролем»}
\label{fig:idef0}
\end{figure}

\subsection{Диаграмма BPMN процесса AS IS}

\begin{figure}[H]
\centering
\begin{tikzpicture}[scale=0.82, every node/.style={transform shape}, node distance=5mm and 5mm]

% Start
\node[bpmn start] (start) {};
\node[font=\tiny\sffamily, below=2mm of start] {Начало};

% Tasks
\node[bpmn task, right=7mm of start] (t1) {Формирование\\заявки на\\материалы};
\node[bpmn task, right=5mm of t1] (t2) {Согласование\\с руковод-\\ством};

% Gateway
\node[bpmn gateway, right=5mm of t2] (g1) {?};

\node[bpmn task, right=5mm of g1] (t3) {Проведение\\тендера};
\node[bpmn task, right=5mm of t3] (t4) {Выбор\\поставщика};

% Second row — увеличен зазор до 26mm
\node[bpmn task, below=26mm of t4] (t5) {Оформление\\договора};
\node[bpmn task, left=5mm of t5] (t6) {Поставка\\материалов};
\node[bpmn task, left=5mm of t6] (t7) {Контроль\\качества\\на площадке};
\node[bpmn task, left=5mm of t7] (t8) {Составление\\отчёта\\(вручную)};

% Gateway 2
\node[bpmn gateway, left=5mm of t8] (g2) {?};

\node[bpmn task, left=5mm of g2] (t9) {Приёмка и\\подписание\\акта};

% End
\node[bpmn end, left=7mm of t9] (endN) {};
\node[font=\tiny\sffamily, below=2mm of endN] {Конец};

% Arrows
\draw[flow arrow] (start) -- (t1);
\draw[flow arrow] (t1) -- (t2);
\draw[flow arrow] (t2) -- (g1);
\draw[flow arrow] (g1) -- (t3) node[midway, above, font=\tiny] {да};
% «нет» — стрелка опущена на 12mm, подпись рядом
\draw[flow arrow] (g1.south) -- ++(0,-12mm) -| (t1.south);
\node[font=\tiny\sffamily, color=samoletdark, anchor=north west] at ([shift={(2pt,-12mm)}]g1.south) {нет};
\draw[flow arrow] (t3) -- (t4);
\draw[flow arrow] (t4) -- (t5);
\draw[flow arrow] (t5) -- (t6);
\draw[flow arrow] (t6) -- (t7);
\draw[flow arrow] (t7) -- (t8);
\draw[flow arrow] (t8) -- (g2);
\draw[flow arrow] (g2) -- (t9) node[midway, above, font=\tiny] {ок};
% «доработка» — стрелка опущена на 12mm, подпись рядом
\draw[flow arrow] (g2.south) -- ++(0,-12mm) -| (t7.south);
\node[font=\tiny\sffamily, color=samoletdark, anchor=north west] at ([shift={(2pt,-12mm)}]g2.south) {доработка};
\draw[flow arrow] (t9) -- (endN);

% Annotations — с запасом над блоками
\node[font=\tiny\itshape\color{accentred}, above=3mm of t1, text width=22mm, align=center] {ручной ввод\\в Excel/10D};
\node[font=\tiny\itshape\color{accentred}, above=3mm of t8, text width=22mm, align=center] {бумажный\\документ};

\end{tikzpicture}
\caption{BPMN-диаграмма процесса закупок и стройконтроля AS IS}
\label{fig:bpmn_asis}
\end{figure}

Процесс характеризуется высокой долей ручных операций: заявки формируются в Excel или вручную вводятся в 10D, отчёты о контроле качества составляются на бумаге, согласования проходят через цепочку email-переписок. Среднее время полного цикла закупки~--- от 14 до 21 дня, из которых до 40\% уходит на ожидание согласований и ручную обработку документов.

\subsection{Диаграмма потоков данных (DFD)}

\begin{figure}[H]
\centering
\begin{tikzpicture}[
  dfd proc/.style={
    circle, draw=samoletblue, fill=samoletlight,
    minimum size=20mm, align=center, font=\small\sffamily
  },
  dfd ext/.style={
    rectangle, draw=samoletdark, fill=samoletdark!8, thick,
    minimum height=10mm, minimum width=30mm, align=center, font=\small\sffamily
  },
  dfd store/.style={
    rectangle, draw=accentteal, fill=accentteal!8,
    minimum height=8mm, minimum width=34mm, align=center, font=\small\sffamily
  }
]

% External entities
\node[dfd ext] (prr) at (-6, 3) {Прораб\\(площадка)};
\node[dfd ext] (mgmt) at (6, 3) {Руководство};
\node[dfd ext] (suppl) at (6, -3) {Поставщики};

% Processes
\node[dfd proc] (p1) at (0, 3) {1.0\\Обработка\\заявок};
\node[dfd proc] (p2) at (0, -1) {2.0\\Проведение\\тендера};
\node[dfd proc] (p3) at (-6, -3) {3.0\\Контроль\\и отчётность};

% Data stores
\node[dfd store] (ds1) at (-3, 0.5) {D1: База заявок};
\node[dfd store] (ds2) at (3, 0.5) {D2: Реестр поставщиков};

% Flows
\draw[flow arrow] (prr) -- (p1) node[midway, above, font=\scriptsize] {заявка};
\draw[flow arrow] (p1) -- (mgmt) node[midway, above, font=\scriptsize] {на согласование};
\draw[flow arrow] (mgmt) -- (p1) node[midway, below, font=\scriptsize, text=accentgreen] {решение};
\draw[flow arrow] (p1) -- (ds1) node[midway, left, font=\scriptsize] {сохранение};
\draw[flow arrow] (ds1) -- (p2) node[midway, left, font=\scriptsize] {данные};
\draw[flow arrow] (ds2) -- (p2) node[midway, right, font=\scriptsize] {контрагенты};
\draw[flow arrow] (p2) -- (suppl) node[midway, right, font=\scriptsize] {приглашение};
\draw[flow arrow] (suppl) -- (p2) node[midway, left, font=\scriptsize, text=accentgreen] {предложение};
\draw[flow arrow] (p2) -- (p3) node[midway, right, font=\scriptsize] {контракт};
\draw[flow arrow] (p3) -- (prr) node[midway, left, font=\scriptsize] {материалы};
\draw[flow arrow] (p3.north east) -- (mgmt.south west) node[midway, right, font=\scriptsize, sloped] {отчёт};

\end{tikzpicture}
\caption{DFD уровня 0: потоки данных процесса закупок и стройконтроля}
\label{fig:dfd}
\end{figure}

Анализ потоков данных выявляет узкие места: множественные переходы между бумажным и электронным форматами, дублирование ввода данных (прораб вводит заявку, затем та же информация повторно вводится в ERP), отсутствие автоматической маршрутизации согласований.

\subsection{Матрица проблем процесса AS IS}

Для систематизации выявленных проблем составлена матрица, связывающая этапы процесса с типами проблем и их влиянием на бизнес.

\begin{table}[H]
\centering
\caption{Матрица проблем процесса закупок и стройконтроля AS IS}
\label{tab:problem_matrix}
\small
\begin{tabularx}{\textwidth}{|l|X|l|l|}
\hline
\rowcolor{samoletlight}
\textbf{Этап} & \textbf{Проблема} & \textbf{Тип} & \textbf{Влияние} \\
\hline
Формирование заявки & Ручной ввод в Excel/10D & Рутина & Высокое \\
\hline
Согласование & Email-цепочки, отсутствие SLA & Задержки & Критическое \\
\hline
Тендер & Ручной сбор предложений & Рутина & Высокое \\
\hline
Выбор поставщика & Субъективные критерии & Качество & Среднее \\
\hline
Поставка & Отсутствие трекинга & Контроль & Высокое \\
\hline
Контроль качества & Визуальный осмотр, бумажные акты & Точность & Критическое \\
\hline
Формирование отчёта & Ручное заполнение, дублирование & Рутина & Высокое \\
\hline
Приёмка & Бумажное подписание & Скорость & Среднее \\
\hline
\end{tabularx}
\end{table}

Из матрицы видно, что наиболее критичные проблемы сосредоточены на этапах согласования (отсутствие автоматической маршрутизации и SLA) и контроля качества (субъективность визуального осмотра и бумажная документация). Именно эти этапы становятся приоритетными для автоматизации через ИИ-платформу.

\newpage

% ============================================================
%  РАЗДЕЛ 8. НОВЫЙ ИТ-ПРОДУКТ
% ============================================================
\section{Описание нового ИТ-продукта}

\subsection{VISION: ИИ-платформа стройконтроля и автоматизации закупок}

Для решения выявленных проблем предлагается внедрение \textbf{мобильной ИИ-платформы} на базе существующей 10D, которая объединит функции стройконтроля, автоматизации закупок и аналитики в единое мобильное решение для линейного персонала.

Платформа «Самолет.ИИ-Стройконтроль» направлена на:
\begin{itemize}[leftmargin=20pt]
  \item Сокращение ручного ввода данных на 70--80\% через ИИ-распознавание фотоматериалов и OCR документов.
  \item Автоматизацию цикла закупок через ИИ-бота для тендеров и автоматическую генерацию заявок.
  \item Предиктивную аналитику для прогнозирования задержек строительства и оптимизации графиков поставок.
  \item Мобильные дашборды для линейного персонала с доступом ко всем данным 10D.
\end{itemize}

\subsection{Product Vision Board}

\begin{figure}[H]
\centering
\begin{tikzpicture}[
  vbox/.style={
    rectangle, rounded corners=4pt, draw=#1!60, fill=#1!8,
    minimum width=145mm, anchor=north west, font=\small\sffamily
  },
  vlabel/.style={font=\small\sffamily\bfseries, color=#1!80}
]

\def\lx{0.4}       % label x
\def\lgap{-0.4}    % label y offset from box top
\def\tgap{-1.15}   % text y offset from box top

% Vision
\node[vbox=samoletblue, minimum height=18mm] (vis) at (0, 0) {};
\node[vlabel=samoletblue, anchor=west] at (\lx, \lgap) {VISION};
\node[font=\small, anchor=west, text width=135mm] at (\lx, -1.05) {
  Мобильная ИИ-платформа, сокращающая рутину линейного персонала на 30--40\%
};

% Target audience
\node[vbox=accentgreen, minimum height=26mm] (ta) at (0, -2.2) {};
\node[vlabel=accentgreen, anchor=west] at (\lx, -2.2+\lgap) {ЦЕЛЕВАЯ АУДИТОРИЯ};
\node[font=\small, anchor=north west, text width=135mm] at (\lx, -2.2+\tgap) {
  Прорабы, инженеры стройконтроля, менеджеры по закупкам, руководители проектов (500+ пользователей)
};

% Needs
\node[vbox=accentorange, minimum height=26mm] (needs) at (0, -5.2) {};
\node[vlabel=accentorange, anchor=west] at (\lx, -5.2+\lgap) {ПОТРЕБНОСТИ};
\node[font=\small, anchor=north west, text width=135mm] at (\lx, -5.2+\tgap) {
  Быстрый ввод данных с площадки, автоматические отчёты, прозрачность закупок, прогноз задержек
};

% Key features
\node[vbox=accentpurple, minimum height=26mm] (feat) at (0, -8.2) {};
\node[vlabel=accentpurple, anchor=west] at (\lx, -8.2+\lgap) {КЛЮЧЕВЫЕ ФУНКЦИИ};
\node[font=\small, anchor=north west, text width=135mm] at (\lx, -8.2+\tgap) {
  ИИ-фото (дефекты), OCR документов, чат-бот тендеров, мобильный дашборд, push-уведомления
};

% Business goals
\node[vbox=accentred, minimum height=26mm] (biz) at (0, -11.2) {};
\node[vlabel=accentred, anchor=west] at (\lx, -11.2+\lgap) {БИЗНЕС-ЦЕЛИ};
\node[font=\small, anchor=north west, text width=135mm] at (\lx, -11.2+\tgap) {
  Сокращение цикла закупки с 21 до 10 дней, экономия 15\% на операциях, рост качества стройконтроля
};

\end{tikzpicture}
\caption{Product Vision Board: «Самолет.ИИ-Стройконтроль»}
\label{fig:vision}
\end{figure}

\subsection{Service Journey Map}

На рисунке~\ref{fig:sjm} представлен путь линейного сотрудника (прораба) при использовании новой ИИ-платформы.

\begin{figure}[H]
\centering
\begin{tikzpicture}[
  step/.style={
    rectangle, rounded corners=5pt, draw=#1!60, fill=#1!12,
    minimum height=18mm, minimum width=26mm, align=center,
    font=\small\sffamily
  },
  sarrow/.style={-{Stealth[length=3mm]}, thick, color=samoletblue!60}
]

\node[step=samoletblue] (s1) at (0, 0) {Фото\\дефекта\\на площадке};
\node[step=accentpurple] (s2) at (3.5, 0) {ИИ-распо-\\знавание\\и оценка};
\node[step=accentorange] (s3) at (7, 0) {Авто-\\создание\\заявки};
\node[step=accentteal] (s4) at (10.5, 0) {Бот\\проводит\\тендер};
\node[step=accentgreen] (s5) at (14, 0) {Дашборд:\\статус\\поставки};

\draw[sarrow] (s1) -- (s2);
\draw[sarrow] (s2) -- (s3);
\draw[sarrow] (s3) -- (s4);
\draw[sarrow] (s4) -- (s5);

% Emotions / satisfaction
\node[font=\scriptsize\sffamily, color=accentgreen!80, align=center] at (0, -1.6) {Удобно:\\1 фото};
\node[font=\scriptsize\sffamily, color=accentgreen!80, align=center] at (3.5, -1.6) {Быстро:\\5 секунд};
\node[font=\scriptsize\sffamily, color=accentgreen!80, align=center] at (7, -1.6) {Без ручного\\ввода};
\node[font=\scriptsize\sffamily, color=accentgreen!80, align=center] at (10.5, -1.6) {Прозрачный\\выбор};
\node[font=\scriptsize\sffamily, color=accentgreen!80, align=center] at (14, -1.6) {Push-\\уведомление};

% Timeline
\draw[gray!40, dashed] (-1.5, -2.3) -- (15.5, -2.3);
\node[font=\tiny\sffamily\color{gray}] at (0, -2.7) {0 мин};
\node[font=\tiny\sffamily\color{gray}] at (3.5, -2.7) {1 мин};
\node[font=\tiny\sffamily\color{gray}] at (7, -2.7) {2 мин};
\node[font=\tiny\sffamily\color{gray}] at (10.5, -2.7) {1--3 дня};
\node[font=\tiny\sffamily\color{gray}] at (14, -2.7) {реалтайм};

\end{tikzpicture}
\caption{Service Journey Map: путь прораба в ИИ-платформе}
\label{fig:sjm}
\end{figure}

\subsection{MVP (Minimum Viable Product)}

Минимально жизнеспособный продукт включает три ключевых модуля:

\begin{enumerate}[leftmargin=20pt]
  \item \textbf{Модуль ИИ-фотофиксации:} мобильное приложение для фотографирования дефектов и хода строительства с автоматическим распознаванием типа дефекта, оценкой серьёзности и привязкой к объекту в 10D.
  \item \textbf{Модуль OCR-документов:} сканирование накладных, актов, счетов-фактур с автоматическим извлечением данных и загрузкой в ТУРБО ERP. Ожидаемая точность распознавания --- не менее 95\%.
  \item \textbf{Мобильный дашборд:} интерфейс для прораба и инженера с отображением статусов объектов, закупок, графиков работ, KPI в реальном времени с push-уведомлениями о критических событиях.
\end{enumerate}

Планируемый срок разработки MVP --- 3 месяца (Q1 2026), пилотирование на 2--3 стройплощадках --- Q2 2026, масштабирование --- Q3--Q4 2026.

\subsection{Стек технологий MVP}

Для реализации MVP выбран следующий стек технологий, обоснованный требованиями к производительности, масштабируемости и интеграции с существующей инфраструктурой:

\begin{table}[H]
\centering
\caption{Технологический стек MVP}
\label{tab:tech_stack}
\small
\begin{tabularx}{\textwidth}{|l|l|X|}
\hline
\rowcolor{samoletlight}
\textbf{Компонент} & \textbf{Технология} & \textbf{Обоснование} \\
\hline
Мобильное приложение & Flutter & Кросс-платформенность (iOS + Android), быстрая разработка \\
\hline
Backend API & Python (FastAPI) & Интеграция с ML-библиотеками, высокая производительность \\
\hline
ИИ-распознавание & YOLOv8 + fine-tuning & Детекция дефектов на изображениях стройплощадок \\
\hline
OCR-движок & Tesseract + доп. обучение & Распознавание накладных, актов, счетов-фактур \\
\hline
Data Lake & Apache Spark + Delta Lake & Обработка больших объёмов данных из всех юнитов \\
\hline
Бот-тендер & LLM + RAG & Генерация тендерной документации, анализ предложений \\
\hline
Интеграция & Apache Kafka & Асинхронный обмен данными между 10D, ERP и мобильным приложением \\
\hline
\end{tabularx}
\end{table}

Выбор стека ориентирован на совместимость с существующей инфраструктурой (Linux, VMware, SAN) и возможность развёртывания как в on-premise ЦОД, так и в гибридном облаке на следующих этапах. Использование open-source компонентов (Python, Tesseract, Spark) снижает лицензионные затраты и соответствует курсу на импортозамещение.

\newpage

% ============================================================
%  РАЗДЕЛ 9. DIGITAL-ТРАНСФОРМАЦИЯ
% ============================================================
\section{Digital-трансформация деятельности компании}

Внедрение ИИ-платформы стройконтроля запускает волну цифровой трансформации, затрагивающую все уровни ИТ-архитектуры и бизнес-процессы компании. На рисунке~\ref{fig:transform} показана схема трансформации по слоям архитектуры.

\begin{figure}[H]
\centering
\begin{tikzpicture}[
  asbox/.style={
    rectangle, rounded corners=3pt, draw=accentred!60, fill=accentred!8,
    minimum height=12mm, minimum width=50mm, align=center, font=\small\sffamily
  },
  tobox/.style={
    rectangle, rounded corners=3pt, draw=accentgreen!60, fill=accentgreen!8,
    minimum height=12mm, minimum width=50mm, align=center, font=\small\sffamily
  },
  tarrow/.style={-{Stealth[length=3mm]}, very thick, color=samoletblue}
]

% AS IS column
\node[font=\normalsize\sffamily\bfseries, color=accentred] at (-4, 6.5) {AS IS};

\node[asbox] (as1) at (-4, 5) {Ручной ввод заявок};
\node[asbox] (as2) at (-4, 3) {Бумажные отчёты};
\node[asbox] (as3) at (-4, 1) {Email-согласования};
\node[asbox] (as4) at (-4, -1) {25+ разрозненных API};
\node[asbox] (as5) at (-4, -3) {Гибридные ЦОД (on-prem)};

% TO BE column
\node[font=\normalsize\sffamily\bfseries, color=accentgreen] at (4, 6.5) {TO BE};

\node[tobox] (to1) at (4, 5) {ИИ-распознавание\\+ авто-заявки};
\node[tobox] (to2) at (4, 3) {Мобильные дашборды\\реального времени};
\node[tobox] (to3) at (4, 1) {Автоматическая\\маршрутизация};
\node[tobox] (to4) at (4, -1) {Единый API-шлюз\\+ data-lake};
\node[tobox] (to5) at (4, -3) {Гибридное облако\\+ резервирование};

% Transformation arrows
\draw[tarrow] (as1) -- (to1);
\draw[tarrow] (as2) -- (to2);
\draw[tarrow] (as3) -- (to3);
\draw[tarrow] (as4) -- (to4);
\draw[tarrow] (as5) -- (to5);

% Layer labels
\node[font=\scriptsize\sffamily\color{gray}, rotate=90] at (-7.5, 5) {Бизнес-процессы};
\node[font=\scriptsize\sffamily\color{gray}, rotate=90] at (-7.5, 3) {Отчётность};
\node[font=\scriptsize\sffamily\color{gray}, rotate=90] at (-7.5, 1) {Согласования};
\node[font=\scriptsize\sffamily\color{gray}, rotate=90] at (-7.5, -1) {Интеграция};
\node[font=\scriptsize\sffamily\color{gray}, rotate=90] at (-7.5, -3) {Инфраструктура};

\end{tikzpicture}
\caption{Схема цифровой трансформации: AS IS $\to$ TO BE по слоям}
\label{fig:transform}
\end{figure}

\subsection{Влияние на деятельность компании}

Цифровая трансформация затрагивает следующие аспекты деятельности:

\begin{enumerate}[leftmargin=20pt]
  \item \textbf{Операционная эффективность:} сокращение рутины линейного персонала на 30--40\% за счёт автоматизации ввода данных, генерации отчётов и маршрутизации согласований. Цикл закупки сокращается с 21 до 10 дней.

  \item \textbf{Качество строительства:} ИИ-распознавание фотоматериалов и данных дронов позволяет выявлять дефекты на ранних стадиях, снижая стоимость исправлений на 50--60\%.

  \item \textbf{Финансовый контроль:} единый data-lake с интеграцией данных всех юнитов обеспечивает прозрачность затрат и позволяет внедрить ITFM (IT Financial Management) для точного учёта расходов на ИТ по проектам.

  \item \textbf{Предиктивная аналитика:} модели машинного обучения на 2,4~ПБ исторических данных прогнозируют задержки строительства с точностью до 85\%, позволяя заблаговременно корректировать графики.

  \item \textbf{Управление рисками:} автоматический мониторинг контрагентов, кредитных рисков и юридических споров через ИИ-анализ открытых источников, что критично в условиях текущего кризиса.
\end{enumerate}

Совокупный эффект трансформации оценивается в экономию 15\% на операционных расходах (около 1,5--2~млрд~руб. в год) и удвоение скорости циклов строительства и закупок к 2027 году.

\subsection{Этапы цифровой трансформации}

Цифровая трансформация ГК «Самолет» реализуется в три последовательных волны, каждая из которых расширяет степень автоматизации и интеллектуализации бизнес-процессов.

\textbf{Волна 1 (Q1--Q2 2026) --- Фундамент.} Создание централизованной data-платформы (data-lake) путём консолидации данных из всех 9 юридических юнитов. Внедрение первых ИИ-модулей в платформу 10D: автоматическое распознавание дефектов на фотографиях стройплощадок, OCR для накладных и актов. Разработка и пилотирование MVP мобильного приложения для прорабов на 2--3 стройплощадках.

\textbf{Волна 2 (Q3--Q4 2026) --- Масштабирование.} Развёртывание мобильного приложения на всех стройплощадках компании (50+ объектов). Запуск ИИ-бота для автоматизации тендерных процедур: подбор поставщиков по критериям, формирование конкурсной документации, анализ предложений. Внедрение предиктивной аналитики для прогнозирования задержек строительства и оптимизации графиков поставок.

\textbf{Волна 3 (2027+) --- Полная автоматизация.} Внедрение RPA-ботов для сквозной автоматизации рутинных операций (формирование отчётности, сверка документов, расчёт зарплат подрядчиков). Миграция в гибридное облако для обеспечения масштабируемости (дополнительные мощности за 1--4 часа) и катастрофоустойчивости. Запуск предиктивной аналитики второго поколения с использованием данных IoT-датчиков и дронов.

\subsection{Управление изменениями}

Для успешной цифровой трансформации критически важно управление изменениями на уровне персонала. Запланировано обучение 20\% линейного персонала (первая волна) с последующим масштабированием. Ключевые мероприятия:

\begin{itemize}[leftmargin=20pt]
  \item Программа «Цифровой прораб» --- обучение работе с мобильным приложением и ИИ-инструментами (40 часов на сотрудника).
  \item Институт «Цифровых амбассадоров» --- подготовка лидеров изменений на каждой стройплощадке.
  \item Система мотивации --- KPI, привязанные к использованию цифровых инструментов.
  \item Help-desk и чат-бот поддержки для оперативного решения технических вопросов.
\end{itemize}

\newpage

% ============================================================
%  РАЗДЕЛ 10. ИТ-АРХИТЕКТУРА TO BE — МЕТАМОДЕЛЬ
% ============================================================
\section{План ИТ-архитектуры компании (TO BE)}

\subsection{Целевая метамодель ИТ-архитектуры}

Целевая архитектура сохраняет четырёхслойную структуру, но включает принципиально новые элементы: ИИ-модули, мобильное приложение, централизованный data-lake и гибридное облако.

\begin{figure}[H]
\centering
\begin{tikzpicture}[
  lblock/.style={
    rectangle, rounded corners=3pt, draw=#1!60, fill=#1,
    minimum height=11mm, minimum width=28mm, align=center,
    font=\small\sffamily, inner sep=3pt
  },
  newblock/.style={
    rectangle, rounded corners=3pt, draw=accentgreen!80, fill=accentgreen!15,
    minimum height=11mm, minimum width=28mm, align=center,
    font=\small\sffamily\bfseries, inner sep=3pt
  }
]

% --- Бизнес-слой ---
\fill[layerbiz, rounded corners=4pt] (-0.5, 7.8) rectangle (17.5, 10.2);
\node[font=\small\sffamily\bfseries, anchor=north west, color=accentorange!80] at (-0.2, 10.1) {Бизнес-слой};

\node[lblock=layerbiz] (b1) at (2, 8.8) {Продажи\\жилья};
\node[lblock=layerbiz] (b2) at (5.5, 8.8) {Строительство\\объектов};
\node[newblock] (b3) at (9, 8.8) {ИИ-закупки\\(автотендер)};
\node[lblock=layerbiz] (b4) at (12.5, 8.8) {Управление\\персоналом};
\node[newblock] (b5) at (16, 8.8) {Предиктивная\\аналитика};

% --- Прикладной слой ---
\fill[layerapp, rounded corners=4pt] (-0.5, 4.8) rectangle (17.5, 7.2);
\node[font=\small\sffamily\bfseries, anchor=north west, color=samoletblue!80] at (-0.2, 7.1) {Прикладной слой};

\node[lblock=layerapp] (a1) at (2, 5.8) {CRM-система};
\node[lblock=layerapp] (a2) at (5.5, 5.8) {10D + ИИ-модули\\(стройконтроль)};
\node[lblock=layerapp] (a3) at (9, 5.8) {ТУРБО ERP\\+ бот-тендер};
\node[newblock] (a4) at (12.5, 5.8) {Мобильное\\приложение};
\node[lblock=layerapp] (a5) at (16, 5.8) {HR + ИИ\\подбор кадров};

% --- Слой данных ---
\fill[layerdata, rounded corners=4pt] (-0.5, 1.8) rectangle (17.5, 4.2);
\node[font=\small\sffamily\bfseries, anchor=north west, color=accentgreen!80] at (-0.2, 4.1) {Слой данных};

\node[lblock=layerdata] (d1) at (2, 2.8) {Клиентская\\база};
\node[lblock=layerdata] (d2) at (5.5, 2.8) {Данные\\стройплощадок};
\node[newblock] (d3) at (9, 2.8) {Data Lake\\(централиз.)};
\node[lblock=layerdata] (d4) at (12.5, 2.8) {Финансовые\\данные};
\node[newblock] (d5) at (16, 2.8) {ИИ-модели\\(ML Pipeline)};

% --- Инфраструктурный слой ---
\fill[layerinfra, rounded corners=4pt] (-0.5, -1.2) rectangle (17.5, 1.2);
\node[font=\small\sffamily\bfseries, anchor=north west, color=gray!70] at (-0.2, 1.1) {Инфраструктурный слой};

\node[lblock=layerinfra] (i1) at (2, -0.1) {ЦОД~1\\(основной)};
\node[lblock=layerinfra] (i2) at (5.5, -0.1) {ЦОД~2\\(резервный)};
\node[newblock] (i3) at (9, -0.1) {Гибридное\\облако};
\node[newblock] (i4) at (12.5, -0.1) {Единый\\API-шлюз};
\node[lblock=layerinfra] (i5) at (16, -0.1) {Сеть +\\мониторинг};

% Вертикальные связи (выборочные)
\foreach \x in {1,2,4} {
  \draw[gray!40, thick, -{Stealth[length=2mm]}] (b\x.south) -- (a\x.north);
  \draw[gray!40, thick, -{Stealth[length=2mm]}] (a\x.south) -- (d\x.north);
}
\draw[accentgreen!60, thick, -{Stealth[length=2mm]}] (b3.south) -- (a3.north);
\draw[accentgreen!60, thick, -{Stealth[length=2mm]}] (a3.south) -- (d3.north);
\draw[accentgreen!60, thick, -{Stealth[length=2mm]}] (b5.south) -- (a5.north);
\draw[accentgreen!60, thick, -{Stealth[length=2mm]}] (a5.south) -- (d5.north);
\draw[accentgreen!60, thick, -{Stealth[length=2mm]}] (a4.south) -- (d4.north);

% Горизонтальные связи
\draw[samoletblue!30, thick, <->] (a1.east) -- (a2.west);
\draw[accentgreen!50, thick, <->] (a2.east) -- (a3.west);
\draw[accentgreen!50, thick, <->] (a3.east) -- (a4.west);
\draw[samoletblue!30, thick, <->] (a4.east) -- (a5.west);

% Новые элементы обозначены жирным шрифтом и зелёной обводкой
\node[font=\footnotesize\sffamily\itshape, color=accentgreen!80] at (8.5, -1.8) {Новые элементы выделены зелёной обводкой и жирным шрифтом};

\end{tikzpicture}
\caption{Метамодель целевой ИТ-архитектуры TO BE}
\label{fig:meta_tobe}
\end{figure}

\subsection{Ключевые изменения по слоям}

\begin{table}[H]
\centering
\caption{Изменения ИТ-архитектуры по слоям: AS IS $\to$ TO BE}
\label{tab:layers_change}
\small
\begin{tabularx}{\textwidth}{|l|X|X|}
\hline
\rowcolor{samoletlight}
\textbf{Слой} & \textbf{AS IS} & \textbf{TO BE} \\
\hline
Бизнес & Ручные закупки, отчёты & ИИ-автотендер, предиктивная аналитика \\
\hline
Приложения & 10D + ERP (без мобильности) & 10D + ИИ + мобильное приложение + бот \\
\hline
Данные & 2,4 ПБ в SAN (фрагментировано) & Централизованный data-lake + ML Pipeline \\
\hline
Инфраструктура & 2 ЦОД (Hitachi/Huawei, 99,95\%) & + гибридное облако + единый API-шлюз (99,99\%) \\
\hline
\end{tabularx}
\end{table}

Переход от текущей к целевой архитектуре предусматривает сохранение ядра (10D и ТУРБО ERP) при надстройке нового слоя: мобильное приложение, ИИ-модули и централизованный data-lake. Это позволяет минимизировать риски миграции и обеспечить поэтапное внедрение.

\subsection{Архитектурные принципы целевого состояния}

При проектировании целевой ИТ-архитектуры были использованы следующие принципы:

\begin{enumerate}[leftmargin=20pt]
  \item \textbf{Принцип эволюционности:} сохранение инвестиций в существующие системы (10D, ТУРБО ERP) с постепенным наращиванием функциональности через ИИ-модули и интеграции.
  \item \textbf{Принцип централизации данных:} единый data-lake как «единая версия правды» для всех подразделений и юридических лиц.
  \item \textbf{Принцип мобильности:} доступ ко всем критичным функциям с мобильных устройств для линейного персонала на стройплощадках.
  \item \textbf{Принцип интеллектуализации:} встраивание ИИ на каждом уровне архитектуры --- от распознавания изображений до предиктивной аналитики управленческих решений.
  \item \textbf{Принцип отказоустойчивости:} повышение SLA с 99,95\% до 99,99\% за счёт гибридного облака и географической распределённости.
  \item \textbf{Принцип масштабируемости:} возможность выделения дополнительных мощностей за 1--4 часа (вместо текущих 2--4 недель) через облачные ресурсы.
\end{enumerate}

\subsection{Миграционная модель (MEA)}

Переход от AS IS к TO BE осуществляется по модели <<параллельного запуска>>: новые компоненты развёртываются параллельно с существующими, затем постепенно переключаются потоки данных. Это обеспечивает бесперебойную работу на протяжении всего периода миграции.

Ключевые риски миграции и меры их митигации:

\begin{table}[H]
\centering
\caption{Риски миграции и меры митигации}
\label{tab:risks}
\small
\begin{tabularx}{\textwidth}{|X|l|X|}
\hline
\rowcolor{samoletlight}
\textbf{Риск} & \textbf{Вероятность} & \textbf{Митигация} \\
\hline
Потеря данных при миграции в data-lake & Средняя & Полное резервное копирование, валидация контрольных сумм \\
\hline
Отторжение ИИ-инструментов персоналом & Высокая & Программа обучения «Цифровой прораб», мотивация через KPI \\
\hline
Несовместимость legacy-API & Средняя & Единый API-шлюз с адаптерами для legacy-систем \\
\hline
Превышение бюджета & Средняя & Резерв 15\%, поэтапное финансирование \\
\hline
Сбои в работе 10D при интеграции ИИ & Низкая & Песочница для тестирования, canary-развёртывание \\
\hline
\end{tabularx}
\end{table}

\newpage

% ============================================================
%  РАЗДЕЛ 11. ДИАГРАММА ГАНТА
% ============================================================
\section{Диаграмма Ганта}

План реализации проекта цифровой трансформации разделён на три этапа с горизонтом планирования до 2027 года. Диаграмма Ганта представлена на рисунке~\ref{fig:gantt}.

\begin{figure}[H]
\centering
\begin{ganttchart}[
  x unit=7mm,
  y unit title=8mm,
  y unit chart=7mm,
  title height=1,
  bar height=0.6,
  bar top shift=0.2,
  vgrid={*{3}{draw=none}, dotted},
  hgrid,
  title label font=\small\sffamily\bfseries,
  bar label font=\small\sffamily,
  group label font=\small\sffamily\bfseries,
  milestone label font=\small\sffamily\itshape,
  bar/.append style={fill=samoletblue!60, draw=samoletblue},
  bar incomplete/.append style={fill=samoletblue!20},
  group/.append style={draw=samoletdark, fill=samoletdark!30},
  milestone/.append style={fill=accentred, draw=accentred!80},
  today=4,
  today label=Сейчас,
  today label font=\tiny\sffamily\itshape,
  today rule/.style={draw=accentred, thick, dashed}
]{1}{18}

\gantttitle{2026}{12} \gantttitle{2027}{6} \\
\gantttitlelist{1,...,12}{1} \gantttitlelist{1,...,6}{1} \\

% Этап 1
\ganttgroup{Этап 1: Краткосрочный}{1}{6} \\
\ganttbar{Централизация data-lake}{1}{4} \\
\ganttbar{Интеграция ИИ в 10D}{2}{5} \\
\ganttbar{Миграция 50\% legacy}{3}{6} \\
\ganttbar{Разработка MVP (моб. прил.)}{1}{3} \\
\ganttmilestone{Запуск MVP пилот}{3} \\

% Этап 2
\ganttgroup{Этап 2: Среднесрочный}{7}{12} \\
\ganttbar{Внедрение ITFM}{7}{9} \\
\ganttbar{Масштабирование моб. прил.}{7}{10} \\
\ganttbar{ИИ-прогнозы задержек}{8}{11} \\
\ganttbar{Бот-тендер}{9}{12} \\
\ganttmilestone{Полное развёртывание}{12} \\

% Этап 3
\ganttgroup{Этап 3: Долгосрочный}{13}{18} \\
\ganttbar{Полный ИИ-стек (RPA)}{13}{16} \\
\ganttbar{Гибридное облако}{13}{17} \\
\ganttbar{Предиктивная аналитика v2}{15}{18} \\
\ganttmilestone{Целевая архитектура}{18}

% Links
\ganttlink{elem4}{elem5}
\ganttlink{elem5}{elem8}
\ganttlink{elem10}{elem11}
\ganttlink{elem11}{elem12}

\end{ganttchart}
\caption{Диаграмма Ганта: план реализации проекта (2026--2027)}
\label{fig:gantt}
\end{figure}

\subsection{Описание этапов}

\textbf{Этап 1 (Q1--Q2 2026):} основной фокус на создании фундамента --- централизация data-lake для консолидации данных из всех юнитов, интеграция первых ИИ-модулей в платформу 10D для стройконтроля, миграция 50\% legacy-систем в современную SAN-архитектуру, разработка и пилотирование MVP мобильного приложения. Ожидаемый эффект: снижение рутины на 25\%, рост доступности до 99,99\%.

\textbf{Этап 2 (Q3--Q4 2026):} масштабирование --- внедрение ITFM для точного учёта ИТ-затрат по проектам, развёртывание мобильного приложения на всех стройплощадках, запуск ИИ-прогнозов задержек и бота для автоматизации тендеров. Ожидаемый эффект: экономия 15\% на операциях, автоматизация 70\% закупочных процедур.

\textbf{Этап 3 (2027+):} завершение трансформации --- полный ИИ-стек с RPA для сквозной автоматизации, миграция в гибридное облако с масштабированием мощностей за 1--4 часа, запуск предиктивной аналитики второго поколения. Ожидаемый эффект: удвоение скорости циклов строительства и закупок.

\newpage

% ============================================================
%  ROADMAP ПРОЕКТА
% ============================================================
\section{Roadmap проекта}

Roadmap визуализирует стратегические потоки работ по кварталам, связывая бизнес-цели, продуктовые фичи и технологические инициативы в единую временну\'{ю} карту.

\begin{figure}[H]
\centering
\begin{tikzpicture}[scale=0.78, every node/.style={transform shape}]

\tikzset{
  qheader/.style={
    rectangle, rounded corners=2pt, draw=samoletblue!70, fill=samoletblue!15,
    minimum height=9mm, minimum width=28mm, align=center,
    font=\small\sffamily\bfseries
  },
  rlabel/.style={
    rectangle, rounded corners=2pt, draw=#1!60, fill=#1!10,
    minimum height=8mm, minimum width=22mm, align=center,
    font=\scriptsize\sffamily\bfseries, text=#1!80
  },
  rbar/.style={
    rectangle, rounded corners=3pt, draw=#1!50, fill=#1!18,
    minimum height=8mm, align=center, font=\scriptsize\sffamily,
    inner xsep=3pt
  }
}

% Quarter headers
\node[qheader] (q1) at (0, 0) {Q1 2026};
\node[qheader] (q2) at (3.4, 0) {Q2 2026};
\node[qheader] (q3) at (6.8, 0) {Q3 2026};
\node[qheader] (q4) at (10.2, 0) {Q4 2026};
\node[qheader] (q5) at (13.6, 0) {2027+};

% --- Row labels ---
\node[rlabel=accentorange, anchor=east] at (-1.8, -1.4) {Бизнес-цели};
\node[rlabel=samoletblue, anchor=east] at (-1.8, -3.0) {Продукт};
\node[rlabel=accentpurple, anchor=east] at (-1.8, -4.6) {Технологии};
\node[rlabel=accentgreen, anchor=east] at (-1.8, -6.2) {Персонал};

% --- Row 1: Бизнес-цели ---
\node[rbar=accentorange, minimum width=28mm] at (1.7, -1.4) {Снижение рутины\\на 25\%};
\node[rbar=accentorange, minimum width=28mm] at (5.1, -1.4) {Пилот на\\3 площадках};
\node[rbar=accentorange, minimum width=28mm] at (8.5, -1.4) {Экономия 15\%\\на операциях};
\node[rbar=accentorange, minimum width=28mm] at (11.9, -1.4) {Автоматизация\\70\% закупок};

% --- Row 2: Продукт ---
\node[rbar=samoletblue, minimum width=60mm] at (1.7, -3.0) {MVP: ИИ-фото + OCR + дашборд};
\node[rbar=samoletblue, minimum width=28mm] at (6.8, -3.0) {Мобильное прил.\\v2.0};
\node[rbar=samoletblue, minimum width=28mm] at (10.2, -3.0) {Бот-тендер\\v1.0};
\node[rbar=samoletblue, minimum width=22mm] at (13.6, -3.0) {Предикт.\\аналитика v2};

% --- Row 3: Технологии ---
\node[rbar=accentpurple, minimum width=28mm] at (0, -4.6) {Централизация\\data-lake};
\node[rbar=accentpurple, minimum width=28mm] at (3.4, -4.6) {Миграция 50\%\\legacy};
\node[rbar=accentpurple, minimum width=28mm] at (6.8, -4.6) {ITFM\\внедрение};
\node[rbar=accentpurple, minimum width=28mm] at (10.2, -4.6) {Единый\\API-шлюз};
\node[rbar=accentpurple, minimum width=22mm] at (13.6, -4.6) {Гибридное\\облако};

% --- Row 4: Персонал ---
\node[rbar=accentgreen, minimum width=60mm] at (1.7, -6.2) {Программа «Цифровой прораб» (20\%)};
\node[rbar=accentgreen, minimum width=28mm] at (6.8, -6.2) {Масштабирование\\обучения};
\node[rbar=accentgreen, minimum width=50mm] at (11.9, -6.2) {Цифровые амбассадоры\\на всех площадках};

% Horizontal separator lines
\foreach \y in {-0.7, -2.2, -3.8, -5.4} {
  \draw[gray!20] (-3, \y) -- (15.2, \y);
}

\end{tikzpicture}
\caption{Roadmap проекта цифровой трансформации ГК «Самолет»}
\label{fig:roadmap}
\end{figure}

Roadmap отражает четыре стратегических потока:

\begin{itemize}[leftmargin=20pt]
  \item \textbf{Бизнес-цели} --- квартальные целевые показатели эффективности, привязанные к KPI руководства.
  \item \textbf{Продукт} --- поэтапное развитие ИИ-платформы от MVP до полнофункциональной системы с предиктивной аналитикой.
  \item \textbf{Технологии} --- инфраструктурные инициативы: data-lake, миграция legacy, API-шлюз, гибридное облако.
  \item \textbf{Персонал} --- программы обучения и управления изменениями, обеспечивающие принятие новых инструментов.
\end{itemize}

\newpage

% ============================================================
%  PRODUCT BACKLOG
% ============================================================
\section{Product Backlog}

Product Backlog систематизирует все задачи проекта с приоритезацией по методу MoSCoW (Must / Should / Could / Won't). Backlog является живым документом, который обновляется по итогам каждого спринта.

\subsection{Backlog: ИИ-платформа «Самолет.ИИ-Стройконтроль»}

\begin{table}[H]
\centering
\caption{Product Backlog с приоритизацией MoSCoW}
\label{tab:backlog}
\small
\begin{tabularx}{\textwidth}{|c|l|X|l|l|}
\hline
\rowcolor{samoletlight}
\textbf{ID} & \textbf{Приоритет} & \textbf{User Story / Задача} & \textbf{Оценка} & \textbf{Спринт} \\
\hline
\multicolumn{5}{|l|}{\cellcolor{accentred!12}\textbf{Must Have --- критически необходимо}} \\
\hline
B-01 & Must & Как прораб, я хочу сфотографировать дефект и получить автоматическую классификацию & 13 SP & S1--S2 \\
\hline
B-02 & Must & Как прораб, я хочу сканировать накладную и автоматически загрузить данные в ERP & 8 SP & S2--S3 \\
\hline
B-03 & Must & Как прораб, я хочу видеть дашборд с KPI объекта на мобильном телефоне & 8 SP & S1--S3 \\
\hline
B-04 & Must & Как менеджер, я хочу чтобы заявка на закупку формировалась автоматически из дефекта & 13 SP & S3--S4 \\
\hline
B-05 & Must & Как архитектор, я хочу централизованный data-lake для консолидации данных юнитов & 21 SP & S1--S4 \\
\hline
\multicolumn{5}{|l|}{\cellcolor{accentorange!12}\textbf{Should Have --- важно}} \\
\hline
B-06 & Should & Как менеджер по закупкам, я хочу бота, который подбирает поставщиков по критериям & 13 SP & S5--S6 \\
\hline
B-07 & Should & Как руководитель, я хочу получать push-уведомления о критических отклонениях & 5 SP & S4 \\
\hline
B-08 & Should & Как фин. директор, я хочу видеть ИТ-затраты по проектам через ITFM & 13 SP & S5--S7 \\
\hline
B-09 & Should & Как прораб, я хочу электронную приёмку с ЭЦП вместо бумажных актов & 8 SP & S6--S7 \\
\hline
B-10 & Should & Как инженер, я хочу чтобы данные с дронов автоматически попадали в 10D & 13 SP & S5--S7 \\
\hline
\multicolumn{5}{|l|}{\cellcolor{accentgreen!12}\textbf{Could Have --- желательно}} \\
\hline
B-11 & Could & Как руководитель, я хочу предиктивный прогноз задержек с точностью 85\% & 21 SP & S8--S10 \\
\hline
B-12 & Could & Как HR, я хочу ИИ-помощника для подбора кадров на стройплощадки & 13 SP & S9--S10 \\
\hline
B-13 & Could & Как менеджер, я хочу RPA-ботов для автоматической сверки документов & 13 SP & S10--S12 \\
\hline
B-14 & Could & Как архитектор, я хочу автомасштабирование через гибридное облако & 21 SP & S10--S12 \\
\hline
\multicolumn{5}{|l|}{\cellcolor{gray!10}\textbf{Won't Have (в этом релизе) --- отложено}} \\
\hline
B-15 & Won't & Полная замена ТУРБО ERP на собственную систему & --- & Бэклог \\
\hline
B-16 & Won't & AR-визуализация BIM-модели на площадке & --- & Бэклог \\
\hline
\end{tabularx}
\end{table}

\subsection{Распределение Backlog по спринтам}

Проект реализуется двухнедельными спринтами (S1--S12). На рисунке~\ref{fig:backlog_flow} показано распределение задач по спринтам и потокам.

\begin{figure}[H]
\centering
\begin{tikzpicture}[scale=0.78, every node/.style={transform shape}]

\tikzset{
  sbox/.style={
    rectangle, rounded corners=3pt, draw=#1!60, fill=#1!12,
    minimum height=9mm, minimum width=22mm, align=center,
    font=\scriptsize\sffamily
  },
  shead/.style={
    rectangle, rounded corners=2pt, draw=samoletblue!50, fill=samoletblue!8,
    minimum height=7mm, minimum width=18mm, align=center,
    font=\scriptsize\sffamily\bfseries
  }
}

% Sprint headers
\foreach \s/\x in {S1/0, S2/2, S3/4, S4/6, S5/8, S6/10, S7/12, S8/14, S9/16} {
  \node[shead] at (\x, 0) {\s};
}

% Must Have row
\node[font=\scriptsize\sffamily\bfseries, color=accentred!70, anchor=east] at (-1.5, -1) {Must};
\node[sbox=accentred, minimum width=42mm] at (1, -1) {B-01: ИИ-фото};
\node[sbox=accentred, minimum width=42mm] at (5, -1) {B-03: Дашборд};
\node[sbox=accentred, minimum width=22mm] at (3, -1.9) {B-02: OCR};
\node[sbox=accentred, minimum width=42mm] at (7, -1.9) {B-04: Авто-заявка};
\node[sbox=accentred, minimum width=62mm] at (3, -2.8) {B-05: Data-lake};

% Should Have row
\node[font=\scriptsize\sffamily\bfseries, color=accentorange!70, anchor=east] at (-1.5, -4) {Should};
\node[sbox=accentorange] at (6, -4) {B-07: Push};
\node[sbox=accentorange, minimum width=42mm] at (9, -4) {B-06: Бот-тендер};
\node[sbox=accentorange, minimum width=62mm] at (10, -4.9) {B-08: ITFM};
\node[sbox=accentorange, minimum width=42mm] at (11, -5.8) {B-09: ЭЦП};

% Could Have row
\node[font=\scriptsize\sffamily\bfseries, color=accentgreen!70, anchor=east] at (-1.5, -7) {Could};
\node[sbox=accentgreen, minimum width=42mm] at (15, -7) {B-11: Предиктив.};
\node[sbox=accentgreen, minimum width=42mm] at (15, -7.9) {B-12: ИИ HR};

% Separator lines
\draw[gray!20] (-2.5, -3.4) -- (17, -3.4);
\draw[gray!20] (-2.5, -6.4) -- (17, -6.4);

% Phase labels at bottom
\draw[decorate, decoration={brace, amplitude=4pt, mirror}, accentred!60, thick] (-1, -3.1) -- (7, -3.1) node[midway, below=5pt, font=\scriptsize\sffamily\color{accentred!70}] {Этап 1: MVP};
\draw[decorate, decoration={brace, amplitude=4pt, mirror}, accentorange!60, thick] (7.5, -6.1) -- (13, -6.1) node[midway, below=5pt, font=\scriptsize\sffamily\color{accentorange!70}] {Этап 2: Масштабирование};
\draw[decorate, decoration={brace, amplitude=4pt, mirror}, accentgreen!60, thick] (13.5, -8.2) -- (17, -8.2) node[midway, below=5pt, font=\scriptsize\sffamily\color{accentgreen!70}] {Этап 3};

\end{tikzpicture}
\caption{Распределение Product Backlog по спринтам и приоритетам}
\label{fig:backlog_flow}
\end{figure}

\subsection{Метрики Backlog}

\begin{table}[H]
\centering
\caption{Сводные метрики Product Backlog}
\label{tab:backlog_metrics}
\begin{tabular}{lr}
\toprule
\textbf{Метрика} & \textbf{Значение} \\
\midrule
Всего User Stories & 16 \\
Must Have & 5 (31\%) \\
Should Have & 5 (31\%) \\
Could Have & 4 (25\%) \\
Won't Have & 2 (13\%) \\
\midrule
Общая оценка (без Won't) & 183 Story Points \\
Velocity (план) & 18--22 SP / спринт \\
Спринтов в проекте & 12 (24 недели) \\
Размер команды & 8--10 человек \\
\bottomrule
\end{tabular}
\end{table}

Backlog управляется Product Owner из ИТ-блока (команда Алексея Семёнова) и пересматривается на Backlog Refinement каждые 2 недели. Приоритеты корректируются исходя из обратной связи пилотных площадок и изменения бизнес-контекста.

\newpage

% ============================================================
%  РАЗДЕЛ 12. СМЕТА
% ============================================================
\section{Смета проекта}

\subsection{Структура затрат}

\begin{table}[H]
\centering
\caption{Смета проекта цифровой трансформации ИТ-архитектуры, млн руб.}
\label{tab:budget}
\small
\begin{tabularx}{\textwidth}{|l|X|r|r|r|r|}
\hline
\rowcolor{samoletlight}
\textbf{№} & \textbf{Статья затрат} & \textbf{Q1--Q2} & \textbf{Q3--Q4} & \textbf{2027} & \textbf{Итого} \\
\hline
1 & Серверное оборудование и облако & 45 & 30 & 60 & 135 \\
\hline
2 & Лицензии ПО (ИИ, ITFM, аналитика) & 25 & 35 & 20 & 80 \\
\hline
3 & Разработка (мобильное прил., боты, интеграции) & 60 & 80 & 50 & 190 \\
\hline
4 & Обучение персонала (20\% линейки) & 10 & 15 & 10 & 35 \\
\hline
5 & Консалтинг и внедрение & 20 & 25 & 15 & 60 \\
\hline
6 & Техническая поддержка & 10 & 15 & 20 & 45 \\
\hline
7 & Резерв (15\%) & 25 & 30 & 26 & 81 \\
\hline
\rowcolor{samoletlight}
  & \textbf{ИТОГО} & \textbf{195} & \textbf{230} & \textbf{201} & \textbf{626} \\
\hline
\end{tabularx}
\end{table}

\subsection{Распределение бюджета}

\begin{figure}[H]
\centering
\begin{tikzpicture}
\def\R{3.5}

% Разработка 30% (0° -> 109.3°)
\fill[samoletblue!70, draw=white, line width=1.2pt] (0,0) -- (0:\R) arc (0:109.3:\R) -- cycle;
\node[font=\scriptsize\sffamily, align=center, text=white] at (54.65:2.2) {Разработка\\30\%};

% Оборудование 22% (109.3° -> 186.9°)
\fill[accentorange!70, draw=white, line width=1.2pt] (0,0) -- (109.3:\R) arc (109.3:186.9:\R) -- cycle;
\node[font=\scriptsize\sffamily, align=center, text=white] at (148.1:2.2) {Оборудование\\22\%};

% Резерв 13% (186.9° -> 233.5°)
\fill[gray!55, draw=white, line width=1.2pt] (0,0) -- (186.9:\R) arc (186.9:233.5:\R) -- cycle;
\node[font=\scriptsize\sffamily, align=center, text=white] at (210.2:2.2) {Резерв\\13\%};

% Лицензии 13% (233.5° -> 279.6°)
\fill[accentpurple!70, draw=white, line width=1.2pt] (0,0) -- (233.5:\R) arc (233.5:279.6:\R) -- cycle;
\node[font=\scriptsize\sffamily, align=center, text=white] at (256.55:2.2) {Лицензии\\13\%};

% Консалтинг 10% (279.6° -> 314.1°)
\fill[accentteal!70, draw=white, line width=1.2pt] (0,0) -- (279.6:\R) arc (279.6:314.1:\R) -- cycle;
\node[font=\scriptsize\sffamily, align=center, text=white] at (296.85:2.3) {Консалтинг\\10\%};

% Поддержка 7% (314.1° -> 340.0°)
\fill[accentgreen!65, draw=white, line width=1.2pt] (0,0) -- (314.1:\R) arc (314.1:340.0:\R) -- cycle;
\node[font=\scriptsize\sffamily, align=center, text=white] at (327.05:2.5) {Поддержка\\7\%};

% Обучение 5% (340.0° -> 360°)
\fill[accentred!55, draw=white, line width=1.2pt] (0,0) -- (340.0:\R) arc (340.0:360:\R) -- cycle;
\node[font=\scriptsize\sffamily, align=center, text=white] at (350:2.6) {Обучение\\5\%};

\end{tikzpicture}
\caption{Распределение бюджета проекта по статьям затрат}
\label{fig:budget_pie}
\end{figure}

Наибольшую долю бюджета составляет разработка (30\%), включающая создание мобильного приложения, ИИ-модулей, бот-тендера и интеграционных решений. Серверное оборудование и облачная инфраструктура составляют 22\% бюджета, что обусловлено необходимостью расширения мощностей для data-lake и гибридного облака.

\newpage

% ============================================================
%  РАЗДЕЛ 13. КОМПАНИЯ TO BE
% ============================================================
\section{Как компания поменялась (TO BE)}

\subsection{Процесс закупок и стройконтроля TO BE}

\begin{figure}[H]
\centering
\begin{tikzpicture}[scale=0.82, every node/.style={transform shape}, node distance=5mm and 5mm]

% Start
\node[bpmn start] (start) {};
\node[font=\tiny\sffamily, below=2mm of start] {Начало};

% Tasks (automated)
\node[bpmn task, right=7mm of start, fill=bpmngreen] (t1) {ИИ-фото:\\авто-заявка\\на материалы};
\node[bpmn task, right=5mm of t1, fill=bpmngreen] (t2) {Автоматич.\\согласование\\(маршрутиз.)};

% Gateway
\node[bpmn gateway, right=5mm of t2] (g1) {?};

\node[bpmn task, right=5mm of g1, fill=bpmngreen] (t3) {Бот-тендер:\\автоподбор\\поставщика};
\node[bpmn task, right=5mm of t3, fill=bpmngreen] (t4) {Электронный\\контракт};

% Second row
\node[bpmn task, below=18mm of t4, fill=bpmngreen] (t5) {Автоматич.\\трекинг\\поставки};
\node[bpmn task, left=5mm of t5, fill=bpmngreen] (t6) {ИИ-контроль\\качества\\(дрон + фото)};
\node[bpmn task, left=5mm of t6, fill=bpmngreen] (t7) {Авто-генерация\\отчёта\\в 10D};

% Gateway 2
\node[bpmn gateway, left=5mm of t7] (g2) {?};

\node[bpmn task, left=5mm of g2, fill=bpmngreen] (t8) {Цифровая\\приёмка\\(ЭЦП)};

% End
\node[bpmn end, left=7mm of t8] (endN) {};
\node[font=\tiny\sffamily, below=2mm of endN] {Конец};

% Arrows
\draw[flow arrow] (start) -- (t1);
\draw[flow arrow] (t1) -- (t2);
\draw[flow arrow] (t2) -- (g1);
\draw[flow arrow] (g1) -- (t3) node[midway, above, font=\tiny] {авто};
% уточнение — стрелка ниже, текст рядом
\draw[flow arrow] (g1.south) -- ++(0,-10mm) -| (t1.south);
\node[font=\tiny\sffamily, anchor=north west, color=samoletdark] at ([shift={(2pt,-10mm)}]g1.south) {уточнение};
\draw[flow arrow] (t3) -- (t4);
\draw[flow arrow] (t4) -- (t5);
\draw[flow arrow] (t5) -- (t6);
\draw[flow arrow] (t6) -- (t7);
\draw[flow arrow] (t7) -- (g2);
\draw[flow arrow] (g2) -- (t8) node[midway, above, font=\tiny] {ок};
% авто-пересъёмка — стрелка ниже, текст рядом
\draw[flow arrow] (g2.south) -- ++(0,-10mm) -| (t6.south);
\node[font=\tiny\sffamily, anchor=north west, color=samoletdark] at ([shift={(2pt,-10mm)}]g2.south) {авто-пересъёмка};
\draw[flow arrow] (t8) -- (endN);

% Annotation
\node[font=\tiny\itshape\color{accentgreen}, below=3mm of t3, text width=30mm, align=center] {Всё автоматизировано:\\0\% ручного ввода};

\end{tikzpicture}
\caption{BPMN-диаграмма процесса закупок и стройконтроля TO BE}
\label{fig:bpmn_tobe}
\end{figure}

\subsection{Сравнение AS IS vs TO BE}

\begin{table}[H]
\centering
\caption{Сравнение ключевых метрик AS IS и TO BE}
\label{tab:comparison}
\begin{tabularx}{\textwidth}{|X|r|r|r|}
\hline
\rowcolor{samoletlight}
\textbf{Метрика} & \textbf{AS IS} & \textbf{TO BE} & \textbf{Улучшение} \\
\hline
Цикл закупки (дни) & 21 & 10 & $-52\%$ \\
\hline
Ручной ввод данных (\% времени) & 30\% & 5\% & $-83\%$ \\
\hline
Время формирования отчёта & 4 часа & 15 мин & $-94\%$ \\
\hline
Доступность ИТ-систем & 99,95\% & 99,99\% & +0,04 п.п. \\
\hline
Автоматизация закупок & 30\% & 85\% & +55 п.п. \\
\hline
Выявление дефектов (стадия) & Приёмка & Ранняя & экономия 50--60\% \\
\hline
Обработка документов (OCR) & 0\% & 95\% & автоматически \\
\hline
Точность прогноза задержек & н/д & 85\% & новая возможность \\
\hline
\end{tabularx}
\end{table}

Переход к целевому состоянию обеспечивает кардинальное улучшение по всем ключевым метрикам: цикл закупки сокращается более чем вдвое, ручной ввод практически исключается, формирование отчётов ускоряется в 16 раз. Появляются принципиально новые возможности --- предиктивная аналитика и автоматическое распознавание дефектов.

\newpage

% ============================================================
%  РАЗДЕЛ 14. ФИНАНСОВЫЙ РЕЗУЛЬТАТ
% ============================================================
\section{Финансовый результат}

\subsection{Анализ инвестиций и экономии}

\begin{table}[H]
\centering
\caption{Финансовая модель проекта, млн руб.}
\label{tab:roi}
\begin{tabular}{lrrrrr}
\toprule
\textbf{Показатель} & \textbf{2026 H1} & \textbf{2026 H2} & \textbf{2027} & \textbf{2028} & \textbf{Итого} \\
\midrule
Инвестиции (CAPEX) & 195 & 230 & 201 & 0 & 626 \\
Операционная экономия & 50 & 120 & 350 & 500 & 1\,020 \\
Сокращение потерь от брака & 20 & 50 & 100 & 150 & 320 \\
Экономия на персонале & 15 & 40 & 80 & 120 & 255 \\
\midrule
\textbf{Совокупная экономия} & \textbf{85} & \textbf{210} & \textbf{530} & \textbf{770} & \textbf{1\,595} \\
\textbf{Чистый эффект (нарастающий)} & $-110$ & $-130$ & $+199$ & $+969$ & \textbf{+969} \\
\bottomrule
\end{tabular}
\end{table}

\subsection{Графики ROI}

\begin{figure}[H]
\centering
\begin{tikzpicture}
\begin{axis}[
  width=140mm,
  height=65mm,
  ylabel={млн руб.},
  xlabel={Период},
  symbolic x coords={2026 H1, 2026 H2, 2027, 2028},
  xtick=data,
  ymin=-250, ymax=1100,
  legend style={at={(0.02,0.98)}, anchor=north west, font=\small},
  grid=major,
  grid style={gray!20},
  every axis plot/.append style={thick},
  axis lines=left,
  xticklabel style={font=\small\sffamily},
  yticklabel style={font=\small\sffamily}
]

\addplot[color=accentred, mark=square*, mark size=2.5pt] coordinates
  {(2026 H1, 195) (2026 H2, 425) (2027, 626) (2028, 626)};

\addplot[color=accentgreen, mark=triangle*, mark size=2.5pt] coordinates
  {(2026 H1, 85) (2026 H2, 295) (2027, 825) (2028, 1595)};

\addplot[color=samoletblue, mark=*, mark size=2.5pt, dashed] coordinates
  {(2026 H1, -110) (2026 H2, -130) (2027, 199) (2028, 969)};

\legend{Инвестиции (кумулятивно), Экономия (кумулятивно), Чистый эффект}

\draw[accentgreen!50, dashed, thin] (axis cs:2026 H2, 0) -- (axis cs:2027, 0);
\node[font=\tiny\sffamily, color=accentgreen!80] at (axis cs:2027, -60) {Точка окупаемости};

\end{axis}
\end{tikzpicture}
\caption{Кумулятивные инвестиции, экономия и чистый эффект проекта}
\label{fig:roi}
\end{figure}

\subsection{Ключевые выводы}

\begin{itemize}[leftmargin=20pt]
  \item \textbf{Общий объём инвестиций:} 626~млн~руб. за период 2026--2027.
  \item \textbf{Совокупная экономия к 2028:} 1\,595~млн~руб. (операционная экономия + снижение потерь + оптимизация ФОТ).
  \item \textbf{Точка окупаемости:} первое полугодие 2027 года (менее 1,5 лет).
  \item \textbf{ROI к концу 2028:} $969 / 626 \times 100\% \approx 155\%$.
  \item \textbf{Payback period:} $\sim$18 месяцев.
\end{itemize}

Проект экономически обоснован: при относительно умеренных инвестициях (менее 1\% от годовой выручки компании) достигается значительная экономия, а срок окупаемости укладывается в 1,5 года. В условиях текущего кризиса ликвидности это особенно важно --- проект позволяет сократить операционные расходы и повысить эффективность использования ограниченных ресурсов.

\subsection{Анализ чувствительности}

Для оценки устойчивости финансовой модели проведён анализ чувствительности ключевых параметров:

\begin{table}[H]
\centering
\caption{Анализ чувствительности финансовой модели}
\label{tab:sensitivity}
\small
\begin{tabularx}{\textwidth}{|X|r|r|r|}
\hline
\rowcolor{samoletlight}
\textbf{Сценарий} & \textbf{Инвестиции} & \textbf{Экономия к 2028} & \textbf{ROI} \\
\hline
Базовый & 626 млн & 1\,595 млн & 155\% \\
\hline
Перерасход бюджета +20\% & 751 млн & 1\,595 млн & 112\% \\
\hline
Экономия ниже на 30\% & 626 млн & 1\,117 млн & 78\% \\
\hline
Задержка внедрения на 6 мес. & 626 млн & 1\,210 млн & 93\% \\
\hline
Пессимистичный (всё вместе) & 751 млн & 940 млн & 25\% \\
\hline
\end{tabularx}
\end{table}

Даже в пессимистичном сценарии (одновременный перерасход бюджета на 20\%, снижение экономии на 30\% и задержка внедрения на 6 месяцев) проект остаётся экономически целесообразным с положительным ROI 25\%. Это подтверждает устойчивость инвестиционного решения и оправдывает его реализацию даже в условиях высокой неопределённости.

% ============================================================
%  ЗАКЛЮЧЕНИЕ
% ============================================================
\section*{Заключение}
\addcontentsline{toc}{section}{Заключение}

В рамках данного проекта проведён комплексный анализ ИТ-архитектуры ГК «Самолет» в условиях глубокого финансового кризиса 2026 года. Исследование охватило все ключевые аспекты: от анализа внешней среды (PESTLE, SWOT) и организационной структуры до детального моделирования бизнес-процессов и проектирования целевой ИТ-архитектуры.

Основные результаты работы:

\begin{enumerate}[leftmargin=20pt]
  \item \textbf{Диагностика текущего состояния:} выявлены критические узкие места --- фрагментация 25+ интеграций, ручной ввод данных на стройплощадках (до 30\% рабочего времени), отсутствие мобильных инструментов и сквозного ИИ.

  \item \textbf{Проектирование решения:} предложена мобильная ИИ-платформа «Самолет.ИИ-Стройконтроль», включающая ИИ-распознавание дефектов, OCR документов, бот-тендер и мобильные дашборды. Разработан MVP и определён стек технологий.

  \item \textbf{Целевая архитектура TO BE:} спроектирована четырёхслойная метамодель с новыми элементами --- data-lake, гибридное облако, единый API-шлюз, ML Pipeline. Обеспечен переход от 99,95\% до 99,99\% доступности.

  \item \textbf{Экономическое обоснование:} проект требует инвестиций в 626~млн~руб. и окупается за 18 месяцев при ROI 155\%. Даже в пессимистичном сценарии ROI остаётся положительным (25\%).
\end{enumerate}

Проект особенно актуален в контексте текущего кризиса компании: цифровая трансформация позволяет сократить операционные расходы на 15\%, ускорить бизнес-циклы вдвое и обеспечить прозрачность управления --- критически важные факторы для восстановления доверия инвесторов и дольщиков.

\newpage

% ============================================================
%  ПРИЛОЖЕНИЯ
% ============================================================
\appendix
\section{Приложение А. Исходные данные для PESTLE-анализа}

\begin{longtable}{|p{25mm}|p{55mm}|p{60mm}|}
\hline
\rowcolor{samoletlight}
\textbf{Фактор} & \textbf{Описание} & \textbf{Источник} \\
\hline
\endfirsthead
\hline
\rowcolor{samoletlight}
\textbf{Фактор} & \textbf{Описание} & \textbf{Источник} \\
\hline
\endhead
Political & Изъятие ЖК «Квартал Марьино» по иску Генпрокуратуры, связанному с Rietumu Banka & monocle.ru (16.02.2026) \\
\hline
Economic & Ключевая ставка ЦБ РФ 21\%, падение спроса на жильё 30--40\% & smart-lab.ru \\
\hline
Social & Иски дольщиков по ЖК «Новое Колпино» и «Курортный квартал» & vedomosti.ru (22.08.2025) \\
\hline
Technological & Централизация ИТ-блока, платформа 10D, 25+ интеграций & ict-online.ru, comnews.ru \\
\hline
Legal & 214-ФЗ, расследования по мошенничеству (закрыты прокуратурой) & repost.news \\
\hline
Environmental & Требования к энергоэффективности зданий класса A & стандарты РФ \\
\hline
\end{longtable}

\section{Приложение Б. Сравнительный анализ ИТ-решений}

\begin{longtable}{|p{30mm}|p{35mm}|p{35mm}|p{35mm}|}
\hline
\rowcolor{samoletlight}
\textbf{Характеристика} & \textbf{10D (текущая)} & \textbf{1С:ERP + BIM} & \textbf{SAP S/4HANA} \\
\hline
\endfirsthead
\hline
\rowcolor{samoletlight}
\textbf{Характеристика} & \textbf{10D (текущая)} & \textbf{1С:ERP + BIM} & \textbf{SAP S/4HANA} \\
\hline
\endhead
Специализация & Девелопмент & Универсальная & Универсальная \\
\hline
Стоимость внедрения & Уже внедрена & 50--100 млн руб. & 200--500 млн руб. \\
\hline
ИИ-возможности & Начальные & Ограниченные & Развитые \\
\hline
Мобильность & Низкая & Средняя & Высокая \\
\hline
Риск миграции & Минимальный & Высокий (1--2 года) & Очень высокий (2--3 года) \\
\hline
Импортозамещение & Российская & Российская & Иностранная (риски) \\
\hline
\textbf{Решение} & \multicolumn{3}{l|}{\textbf{Модернизация 10D + ИИ-надстройка (оптимальный баланс)} } \\
\hline
\end{longtable}

\section{Приложение В. Календарный план работ (роадмап)}

\begin{longtable}{|c|p{40mm}|p{30mm}|p{25mm}|p{30mm}|}
\hline
\rowcolor{samoletlight}
\textbf{№} & \textbf{Задача} & \textbf{Ответственный} & \textbf{Срок} & \textbf{Результат} \\
\hline
\endfirsthead
\hline
\rowcolor{samoletlight}
\textbf{№} & \textbf{Задача} & \textbf{Ответственный} & \textbf{Срок} & \textbf{Результат} \\
\hline
\endhead
1 & Аудит текущей ИТ-архитектуры & ИТ-блок (Семёнов А.) & Фев 2026 & Отчёт AS IS \\
\hline
2 & Проектирование data-lake & Архитектор данных & Фев--мар 2026 & Техпроект \\
\hline
3 & Разработка MVP моб. приложения & Команда разработки & Фев--апр 2026 & MVP v1.0 \\
\hline
4 & Обучение ИИ-модели (дефекты) & Data Science & Мар--апр 2026 & Модель YOLOv8 \\
\hline
5 & Пилот на 3 площадках & Проектный офис & Апр--май 2026 & Пилот-отчёт \\
\hline
6 & Миграция 50\% legacy в SAN & Инфраструктура & Мар--июн 2026 & Миграция \\
\hline
7 & Внедрение ITFM & Финансовый блок & Июл--сен 2026 & ITFM-система \\
\hline
8 & Масштабирование на все площадки & Проектный офис & Июл--окт 2026 & Массовый запуск \\
\hline
9 & Разработка бот-тендера & Команда разработки & Сен--дек 2026 & Бот v1.0 \\
\hline
10 & ИИ-прогнозы задержек & Data Science & Авг--ноя 2026 & ML-модель \\
\hline
11 & Внедрение RPA & Автоматизация & Янв--апр 2027 & RPA-боты \\
\hline
12 & Миграция в гибридное облако & Инфраструктура & Янв--май 2027 & Облачная среда \\
\hline
13 & Предиктивная аналитика v2 & Data Science & Мар--июн 2027 & ML v2.0 \\
\hline
14 & Финальная оценка эффективности & Проектный офис & Июн 2027 & Итоговый отчёт \\
\hline
\end{longtable}

\section{Приложение Г. Детальный расчёт экономии}

\begin{table}[H]
\centering
\small
\begin{tabular}{|p{50mm}|r|r|p{40mm}|}
\hline
\rowcolor{samoletlight}
\textbf{Статья экономии} & \textbf{Текущие затраты, млн/год} & \textbf{Экономия, \%} & \textbf{Обоснование} \\
\hline
Ручной ввод данных (ФОТ) & 180 & 35\% & ИИ + OCR \\
\hline
Потери от брака и переделок & 400 & 25\% & Раннее выявление дефектов \\
\hline
Избыточные закупки & 300 & 15\% & Оптимизация тендеров \\
\hline
Штрафы за просрочки & 200 & 20\% & Предиктивная аналитика \\
\hline
Поддержка legacy-систем & 100 & 40\% & Миграция и консолидация \\
\hline
\rowcolor{samoletlight}
\textbf{ИТОГО} & \textbf{1\,180} & --- & \textbf{$\sim$500 млн/год} \\
\hline
\end{tabular}
\end{table}

\newpage

% ============================================================
%  СПИСОК ИСТОЧНИКОВ
% ============================================================
\renewcommand{\refname}{Список источников}
\begin{thebibliography}{99}

\bibitem{monocle2026} Просадка «Самолета» // Monocle. --- 2026. --- 16 февр. --- URL: \url{https://monocle.ru/monocle/2026/08/prosadka-samoleta/}

\bibitem{repost2026} Прокуратура прекратила расследование в отношении ГК «Самолет» // Repost.news. --- 2026. --- URL: \url{https://repost.news/news/96502}

\bibitem{vedomosti2025} Топ-менеджерам петербургской структуры «Самолета» грозит до 10 лет // Ведомости. --- 2025. --- 22 авг. --- URL: \url{https://www.vedomosti.ru/realty/articles/2025/08/22/1133405}

\bibitem{smartlab2026} Что происходит с ГК «Самолет»? Анализ причин кризиса // Smart-Lab. --- 2026. --- URL: \url{https://smart-lab.ru/blog/1263596.php}

\bibitem{finbazar2026} Кризис ГК «Самолет»: акции рухнули, иски дольщиков // Финбазар. --- 2026. --- URL: \url{https://finbazar.ru/post/321105}

\bibitem{devmedia2026} Самолет приземлился на дозаправку. Что ждет компанию и её дольщиков // Development Media. --- 2026. --- 5 февр. --- URL: \url{https://developmentmedia.ru/2026/02/05/samolet-prizemlilsya}

\bibitem{ict2025} «Самолет» централизует ИТ-блок для создания единой цифровой платформы // ICT-Online. --- 2025. --- URL: \url{https://ict-online.ru/news/Samolet-tsentralizuyet-IT-blok}

\bibitem{comnews2024} Группа «Самолет» построила отказоустойчивую корпоративную ИТ-инфраструктуру // ComNews. --- 2024. --- 4 апр. --- URL: \url{https://www.comnews.ru/content/232444}

\bibitem{ptmk2024} Группа «Самолет» внедрит единую информационную систему // ПТМК. --- 2024. --- URL: \url{https://www.ptmk.ru/samolet/}

\bibitem{tadviser2025} Проекты на базе ИИ в Группе компаний «Самолет» // TAdviser. --- 2025. --- URL: \url{https://www.tadviser.ru}

\bibitem{samolet_ir} Корпоративное управление группы компаний // ГК «Самолет». --- URL: \url{https://samolet.ru/investors/management/}

\bibitem{smartsmi2025} ГК «Самолет» намерена увеличить свою долю // Smart-SMI. --- 2025. --- URL: \url{https://smart-smi.ru/single/2841642}

\bibitem{wiki_samolet} Группа «Самолет» // Википедия. --- URL: \url{https://ru.wikipedia.org/wiki/Группа_«Самолет»}

\bibitem{ratings2026} ГК «Самолет»: рейтинг с A+.ru до A.ru // Рейтинговое агентство. --- 2026. --- 12 февр. --- URL: \url{https://ratings.ru/ratings/press-releases/Samolet-RA-120226/}

\bibitem{mk2026} «Самолет» попал в зону турбулентности // Московский комсомолец. --- 2026. --- 16 февр. --- URL: \url{https://www.mk.ru/realestate/2026/02/16/samolet}

\end{thebibliography}

\end{document}
